% ============================================================
%  ModNet.Flat.tex
%  ModNet: A Self-Growing Layer-Free Recurrent Network
%  for Universal Computation
%
%  Flat architecture (no modules) with full event logging.
%
%  Compilation:
%    tectonic ModNet.Flat.tex
%
%  Each chapter has its own bibliography file:
%    refs_01_intro.bib, refs_02_related.bib, ...
%  The `chapterbib` package handles per-chapter bibliographies.
% ============================================================

\documentclass[11pt,a4paper]{report}

% ------------------------------------------------------------
%  Encoding & Language
% ------------------------------------------------------------
\usepackage[utf8]{inputenc}
\usepackage[T1]{fontenc}
\usepackage[english]{babel}

% ------------------------------------------------------------
%  Unicode character fallbacks (for T1 fonts)
% ------------------------------------------------------------
\usepackage{newunicodechar}
\newunicodechar{—}{---}
\newunicodechar{–}{--}
\newunicodechar{−}{\ensuremath{-}}
\newunicodechar{…}{\ldots}
\newunicodechar{“}{``}
\newunicodechar{”}{''}
\newunicodechar{‘}{`}
\newunicodechar{’}{'}
\newunicodechar{≈}{\ensuremath{\approx}}
\newunicodechar{≥}{\ensuremath{\geq}}
\newunicodechar{≤}{\ensuremath{\leq}}
\newunicodechar{×}{\ensuremath{\times}}
\newunicodechar{÷}{\ensuremath{\div}}
\newunicodechar{±}{\ensuremath{\pm}}
\newunicodechar{⊕}{\ensuremath{\oplus}}
\newunicodechar{⊗}{\ensuremath{\otimes}}
\newunicodechar{→}{\ensuremath{\rightarrow}}
\newunicodechar{←}{\ensuremath{\leftarrow}}
\newunicodechar{↔}{\ensuremath{\leftrightarrow}}
\newunicodechar{∈}{\ensuremath{\in}}
\newunicodechar{∑}{\ensuremath{\sum}}
\newunicodechar{∏}{\ensuremath{\prod}}
\newunicodechar{√}{\ensuremath{\sqrt{}}}
\newunicodechar{∞}{\ensuremath{\infty}}
\newunicodechar{α}{\ensuremath{\alpha}}
\newunicodechar{β}{\ensuremath{\beta}}
\newunicodechar{γ}{\ensuremath{\gamma}}
\newunicodechar{δ}{\ensuremath{\delta}}
\newunicodechar{ε}{\ensuremath{\varepsilon}}
\newunicodechar{θ}{\ensuremath{\theta}}
\newunicodechar{λ}{\ensuremath{\lambda}}
\newunicodechar{μ}{\ensuremath{\mu}}
\newunicodechar{π}{\ensuremath{\pi}}
\newunicodechar{ρ}{\ensuremath{\rho}}
\newunicodechar{σ}{\ensuremath{\sigma}}
\newunicodechar{τ}{\ensuremath{\tau}}
\newunicodechar{φ}{\ensuremath{\varphi}}
\newunicodechar{ψ}{\ensuremath{\psi}}
\newunicodechar{ω}{\ensuremath{\omega}}
\newunicodechar{Δ}{\ensuremath{\Delta}}
\newunicodechar{Σ}{\ensuremath{\Sigma}}
\newunicodechar{Ω}{\ensuremath{\Omega}}
\newunicodechar{·}{\ensuremath{\cdot}}

% ------------------------------------------------------------
%  Typography & Layout
% ------------------------------------------------------------
\usepackage{lmodern}
\usepackage{microtype}
\usepackage{geometry}
\geometry{margin=1in}
\linespread{1.05}

% ------------------------------------------------------------
%  Math
% ------------------------------------------------------------
\usepackage{amsmath}
\usepackage{amssymb}
\usepackage{amsthm}
\usepackage{mathtools}

% ------------------------------------------------------------
%  Graphics
% ------------------------------------------------------------
\usepackage{graphicx}
\usepackage{xcolor}
\usepackage{subcaption}
\usepackage{float}
\usepackage[section]{placeins}

% ------------------------------------------------------------
%  TikZ & PGFPlots
% ------------------------------------------------------------
\usepackage{tikz}
\usepackage{pgfplots}
\pgfplotsset{compat=1.18}
\usetikzlibrary{arrows.meta, positioning, shapes.geometric, fit, backgrounds}

% ------------------------------------------------------------
%  Tables
% ------------------------------------------------------------
\usepackage{booktabs}
\usepackage{multirow}
\usepackage{array}
\usepackage{longtable}
\usepackage{csvsimple}

% ------------------------------------------------------------
%  Algorithms
% ------------------------------------------------------------
\usepackage{algorithm}
\usepackage{algpseudocode}

% ------------------------------------------------------------
%  Code listings
% ------------------------------------------------------------
\usepackage{listings}
\lstset{
  basicstyle=\ttfamily\small,
  breaklines=true,
  frame=single,
  showstringspaces=false,
  keywordstyle=\color{blue!70!black},
  commentstyle=\color{green!50!black},
  stringstyle=\color{red!70!black},
}

% ------------------------------------------------------------
%  Per-chapter bibliographies
%  Each chapter file ends with:
%    \bibliographystyle{plainnat}
%    \bibliography{refs_XX_chaptername}
%  and provides its own .bib file in the same directory.
% ------------------------------------------------------------
% (chapterbib removed by patch v6 — using single bibliography)
\usepackage[numbers,sort&compress]{natbib}

% ------------------------------------------------------------
%  Hyperlinks
% ------------------------------------------------------------
\usepackage[colorlinks=true,
            linkcolor=blue!60!black,
            citecolor=blue!60!black,
            urlcolor=blue!60!black]{hyperref}

% ------------------------------------------------------------
%  Theorem environments
% ------------------------------------------------------------
\newtheorem{theorem}{Theorem}[chapter]
\newtheorem{lemma}[theorem]{Lemma}
\newtheorem{proposition}[theorem]{Proposition}
\newtheorem{corollary}[theorem]{Corollary}
\theoremstyle{definition}
\newtheorem{definition}[theorem]{Definition}
\newtheorem{example}[theorem]{Example}
\theoremstyle{remark}
\newtheorem{remark}[theorem]{Remark}

% ------------------------------------------------------------
%  Custom commands
% ------------------------------------------------------------
\newcommand{\modnet}{\textsc{ModNet}}
\newcommand{\R}{\mathbb{R}}
\newcommand{\B}{\mathbb{B}}
\newcommand{\N}{\mathbb{N}}
\newcommand{\E}{\mathbb{E}}
\newcommand{\vect}[1]{\boldsymbol{#1}}
\newcommand{\mat}[1]{\mathbf{#1}}

% ------------------------------------------------------------
%  Figure paths
%  All benchmark outputs live under results_bench_v2/
%  Each problem has:
%    best_spring.png           - spring-layout graph of the final net
%    snapshot_gen_XXXX.png     - snapshots at selected generations
% ------------------------------------------------------------
\graphicspath{
  {results_bench_v2/}
  {figures/}
}

% ============================================================

% -------------------------------------------------------------
%  Safe includegraphics: placeholder box if the file is missing.
%  \IfFileExists does NOT consult \graphicspath, so we check
%  the full path under results_bench_v2/.
% -------------------------------------------------------------
\newcommand{\safeincludegraphics}[2][]{%
  \IfFileExists{results_bench_v2/#2}%
    {\includegraphics[#1]{#2}}%
    {\fbox{\parbox[c][3cm][c]{0.6\linewidth}{%
       \centering\ttfamily\small
       [missing figure: \detokenize{#2}]}}}%
}

\begin{document}

% ------------------------------------------------------------
%  Title page
% ------------------------------------------------------------
\title{%
  \textbf{ModNet: A Self-Growing Layer-Free Recurrent Network\\[0.3em]
  for Universal Computation}\\[1em]
  \large An evolutionary framework that discovers structure
  from first principles
}

\author{%
  Masoud Azizi\thanks{Corresponding author.
    Email: \texttt{mablue92@gmail.com}} \\[0.5em]
  \small Department of Computer Science \\
  \small Faculty of Mathematical Sciences \\
  \small University of Tabriz \\
  \small Tabriz, Iran
}

\date{\today}

\maketitle

% ------------------------------------------------------------
%  Abstract
% ------------------------------------------------------------
\begin{abstract}
\noindent
We present \modnet{}, a self-growing layer-free recurrent network
whose topology is discovered through evolution rather than imposed
by a designer. The architecture has no predefined layers and no
modular grouping: every node is a general-purpose threshold unit
that may receive input from any other node, including itself, and
the first $n_{\text{out}}$ nodes of the network are designated as
outputs. Starting from a minimal random topology, the network grows,
prunes, and reorganizes itself through a compact evolutionary loop.

We evaluate \modnet{} on a suite of sixteen benchmark problems
spanning Boolean logic, digital arithmetic, sorting networks, and
continuous regression. The network solves fourteen of the sixteen
tasks, including full adder, 2-bit adder, comparator, 4-to-1
multiplexer, and 4-bit sorting network, all with exact $100\%$
accuracy using no more than $27$ nodes. On continuous targets,
\modnet{} achieves root-mean-square errors as low as $0.0143$ on the
$x^2 + y^2$ surface and $0.0254$ on the multiplication task with
only $13$ nodes. We observe that recurrence emerges spontaneously
and scales with task complexity: XOR-2 requires one self-loop, while
Gaussian and multiplication require $10$ to $14$ self-loops,
indicating that temporal memory is a structural necessity, not an
imposed design.

Our results suggest that a fully layer-free architecture driven by
a minimal evolutionary process can rediscover structural principles
— non-linearity, recurrence, sparsity — that are typically imposed
by human designers. We provide a complete open-source implementation
and a reproducible benchmark suite.

\vspace{1em}
\noindent
\textbf{Keywords:} neuroevolution, layer-free architectures,
self-growing networks, recurrent networks, Boolean function
learning, digital circuits, sparse connectivity.

\end{abstract}

% ------------------------------------------------------------
%  Front matter
% ------------------------------------------------------------
\tableofcontents
\listoffigures
\listoftables

% ============================================================
%                       CHAPTERS
% ============================================================
% During writing, uncomment the line below to compile only the
% chapters you are currently working on. This speeds up
% compilation significantly.
%
% \includeonly{
%   chapters/05_results,
% }

% ============================================================
%  Chapter 1: Introduction
%  File: chapters/01_introduction.tex
%  Bibliography: chapters/refs_01_intro.bib
% ============================================================

\chapter{Introduction}
\label{ch:introduction}

\section{Motivation}

The design of artificial neural networks has historically been guided by
human intuition. From the perceptron model of \citet{rosenblatt1958perceptron}
to modern deep architectures \citep{lecun2015deep}, researchers have manually
specified the number of layers, the connectivity patterns, the activation
functions, and the learning rules. This top-down approach has produced
remarkable results, but it has also imposed structural constraints that may
not be optimal for every task.

A central example of such a constraint is the layered architecture itself.
Since the backpropagation algorithm was popularized by
\citet{rumelhart1986learning}, nearly all successful neural networks have
been organized into discrete layers. This design is computationally
convenient---it allows efficient matrix multiplication and gradient
computation---but it also imposes a rigid separation between
``input,'' ``hidden,'' and ``output'' units. In a biological brain, by
contrast, there are no layers. Neurons form a dense, recurrent graph in
which any neuron may connect to any other, and structure emerges through
developmental and evolutionary processes rather than explicit design.

This discrepancy raises a fundamental question:

\begin{quote}
\emph{Can a neural network discover its own structure from first
principles, without a designer specifying layers, connectivity, or even
the number of nodes?}
\end{quote}

The present work answers this question in the affirmative, through a
minimal architecture we call \modnet{}.

\section{Background}

\subsection{The XOR Problem and the Limits of Linear Models}

The limitations of fixed, single-layer architectures were formalized by
\citet{minsky1969perceptrons}, who proved that a single-layer perceptron
cannot compute the XOR function, nor the parity function for any number of
inputs. This result, often referred to as the ``XOR problem,'' is not merely
a technical curiosity. It reflects a deeper mathematical fact: linear
threshold units can only separate linearly separable classes, and parity is
not linearly separable for any $n \geq 2$.

The XOR problem had a profound impact on the field. It contributed to the
first ``AI winter'' and demonstrated that architectural choices---in
particular, the presence or absence of hidden layers---are not merely
implementation details but determine the computational power of the model.
It was not until the rediscovery of backpropagation by
\citet{rumelhart1986learning} that multi-layer networks became trainable,
and the XOR problem was finally solved by adding a single hidden layer.

\subsection{Evolutionary Approaches to Network Design}

An alternative to gradient-based training is to evolve the network topology
itself. The NeuroEvolution of Augmenting Topologies (NEAT) algorithm,
introduced by \citet{stanley2002evolving}, was a landmark contribution in
this direction. NEAT starts with a minimal network and incrementally adds
nodes and connections, using a principled crossover mechanism and
speciation to protect structural innovation. It demonstrated that evolving
topology alongside weights can significantly accelerate learning on
reinforcement learning benchmarks.

NEAT was later extended by HyperNEAT \citep{stanley2009hypercube}, which
uses an indirect encoding called a Compositional Pattern-Producing Network
(CPPN) to generate connectivity patterns with geometric regularity.
HyperNEAT enabled the evolution of much larger networks than direct
encoding, but it still relies on a predefined substrate geometry, typically
a grid of neurons.

A parallel line of work explored growing neural networks with incremental
learning. The Cascade-Correlation architecture of
\citet{fahlman1990cascade} begins with a minimal network and adds hidden
units one by one. The Growing Neural Gas algorithm of
\citet{fritzke1995growing} uses competitive Hebbian learning to build
topological representations incrementally. Both approaches demonstrate the
power of growth, but neither fully addresses the problem of recurrence,
nor the possibility of a fully layer-free topology.

\subsection{Modularity in Neural Networks}

Modularity---the decomposition of a system into cohesive, loosely coupled
components---is a well-established principle in software engineering
\citep{parnas1972modularity}. \citet{alho2024modularity} conducted a
systematic review of $86$ studies on modular neural networks and found that
nearly two-thirds report improvements in task accuracy compared to
monolithic solutions.

In the context of neuroevolution, \citet{khare2005coevolutionary} proposed
a co-evolutionary model for automatic problem decomposition, in which a
population of modules is evolved alongside a population of complete
systems that combine modules. They demonstrated that if a particular task
decomposition is beneficial, it can be evolved through this mechanism.
However, their model requires the designer to specify the number of modules
and the interface between them.

\citet{landassuri2012evolution} investigated the emergence of modularity
in evolved neural networks and found that pure-modular architectures emerge
at low connectivity values. These results suggest that modularity is not
merely a design choice but can be an emergent property of evolutionary
search.

In the present work, we deliberately \emph{abandon} explicit modular
structure. Our architecture has no modules at all. This is a stronger
demonstration of emergence: the properties we observe---recurrence,
sparsity, cross-scale integration---must arise entirely from the
evolutionary process, without any modular scaffold to support them.

\subsection{Sparse Connectivity}

Biological neural networks are highly sparse: the probability that any two
neurons are connected is on the order of $10^{-6}$ \citep{lecun2015deep}.
This sparsity is not accidental; it is a consequence of metabolic
constraints and is believed to be essential for efficient information
processing \citep{chklovskii2004wiring}.

In artificial neural networks, sparsity has been studied primarily as a
compression technique. The Sparse Evolutionary Training (SET) method of
\citet{mocanu2018scalable} evolves an initial sparse topology into a
scale-free topology during training, allowing the training of networks
with over one million neurons on commodity hardware. However, these methods
typically begin with a predefined network architecture and prune or
sparsify it. They do not address the question of how sparsity and
connectivity structure emerge from first principles.

\subsection{Universal Computation}

The theoretical foundations of neural computation were established by
\citet{siegelmann1995computational}, who proved that recurrent neural
networks with rational weights are Turing complete. More recently,
\citet{perez2021turing} showed that transformers and Neural GPUs are also
Turing complete, without requiring external memory. These results
establish that neural networks can, in principle, simulate any computable
function, but they do not address the question of how such networks can be
discovered or evolved.

\section{The Present Work}

In this paper, we present \modnet{}, a self-growing layer-free recurrent
network that addresses the limitations of fixed architectures by
discovering its own structure through evolution. The architecture has
four defining characteristics:

\begin{enumerate}
  \item \textbf{Layer-free.} There are no predefined layers, and no
    modular grouping. Every node is a general-purpose threshold unit that
    may connect to any other node, including itself.

  \item \textbf{Outputs as ordinary nodes.} The first $n_{\text{out}}$
    nodes of the network are designated as outputs; there is no separate
    output layer.

  \item \textbf{Grown and pruned.} The network starts from a small random
    topology and grows through addition of nodes and pruned through removal
    of inactive or unreachable ones. Growth is bidirectional.

  \item \textbf{Gradient-free.} No gradient information is used at any
    stage. All learning is performed by a compact evolutionary algorithm
    operating on the topology and weights simultaneously.
\end{enumerate}

We evaluate \modnet{} on a suite of \textbf{sixteen} benchmark problems
spanning four broad classes: Boolean logic, digital arithmetic,
sorting networks, and continuous regression. Table~\ref{tab:intro-headline}
summarizes the headline results.

\begin{table}[h]
  \centering
  \small
  \caption{Headline results of \modnet{} on the sixteen benchmark
    problems. Boolean tasks are measured in exact accuracy; continuous
    tasks in root-mean-square error (RMSE). Full details are reported
    in Chapter~\ref{ch:results}.}
  \label{tab:intro-headline}
  \begin{tabular}{llrrr}
    \toprule
    \textbf{Task} & \textbf{Type} & \textbf{Result} & \textbf{Nodes} & \textbf{Time (s)} \\
    \midrule
    XOR-2               & Boolean    & \textbf{100.0\%} &  8 &  0.8 \\
    XOR-3               & Boolean    & \textbf{100.0\%} &  9 &  2.0 \\
    Parity-4            & Boolean    & \textbf{100.0\%} & 10 &  1.8 \\
    Parity-6            & Boolean    & 95.3\%           &  8 & 40.6 \\
    Majority-3          & Boolean    & \textbf{100.0\%} & 26 &  0.9 \\
    Majority-5          & Boolean    & \textbf{100.0\%} & 26 & 10.2 \\
    Full Adder          & Boolean    & \textbf{100.0\%} & 26 &  3.0 \\
    2-bit Adder         & Boolean    & \textbf{100.0\%} & 26 &  7.3 \\
    Comparator          & Boolean    & \textbf{100.0\%} & 27 &  3.8 \\
    MUX 4-to-1          & Boolean    & \textbf{100.0\%} & 26 &  4.8 \\
    Sort-4              & Boolean    & \textbf{100.0\%} & 26 &  8.7 \\
    \midrule
    $\sin(2\pi x)$      & Regression & RMSE $= 0.0465$  & 18 & 30.0 \\
    $\sin(4\pi x)$      & Regression & RMSE $= 0.0789$  & 13 & 25.1 \\
    $x^2 + y^2$         & Regression & RMSE $= 0.0143$  & 26 & 21.1 \\
    Gaussian            & Regression & RMSE $= 0.0275$  & 27 & 41.6 \\
    Multiplication      & Regression & RMSE $= 0.0254$  & 13 & 28.5 \\
    \bottomrule
  \end{tabular}
\end{table}

Three observations follow immediately from
Table~\ref{tab:intro-headline}. First, \modnet{} solves \textbf{fourteen
of sixteen} tasks to their success threshold, including all of the digital
arithmetic and sorting tasks. Second, the evolved networks are small: no
task required more than $27$ nodes, and the multiplication task required
only $13$. Third, the same architecture, without any task-specific
modification, handles Boolean logic, digital arithmetic, sorting, and
continuous regression.

\section{Key Findings}

Beyond the raw numbers, we observe several structural properties that
emerge spontaneously from the evolutionary process.

\subsection{Recurrence Emerges, and Scales with Complexity}

Every successfully evolved network contains recurrent self-loops, ranging
from one self-loop (XOR-2) to fourteen (Gaussian). No mechanism in our
algorithm explicitly favors recurrence. Self-loops emerge because the
unrestricted topology allows them, and evolutionary pressure discovers them
to be useful.

Crucially, the number of self-loops scales with the temporal complexity of
the task. XOR-2, a purely combinatorial problem, requires only one
self-loop. Gaussian and multiplication, both continuous, require ten to
fourteen. This suggests that recurrence is a \emph{structural necessity}
for tasks that require temporal integration, and that \modnet{} discovers
this necessity without being told.

\subsection{Sparsity Is the Norm}

Evolved networks use a small fraction of the available connections. Density,
defined as the ratio of active connections to total possible connections
($E / N^2$), ranges from $27\%$ to $67\%$ depending on the task. This
sparsity was not imposed by any regularizer; it emerged from the pruning
mechanism, which removes nodes whose activity is either too low or too
high.

The observed sparsity is consistent with that of biological neural
networks \citep{chklovskii2004wiring}, where the vast majority of
potential synapses are absent.

\subsection{Growth Adapts to Task Complexity}

The initial network has $26$ nodes. Over evolution, networks either
shrank, grew, or stayed approximately the same:

\begin{itemize}
  \item \textbf{Shrinking:} XOR-2 ($26 \to 8$), XOR-3 ($26 \to 9$),
    Parity-4 ($26 \to 10$), Multiplication ($26 \to 13$).
  \item \textbf{Growing:} sin(2$\pi x$) ($26 \to 18$), Gaussian
    ($26 \to 27$).
  \item \textbf{Nearly stable:} Parity-6, Majority-5, Full Adder.
\end{itemize}

Boolean tasks tend to converge to small networks, while continuous tasks
require larger networks. The pattern is consistent with the intuition that
Boolean functions have low algorithmic complexity, while continuous
functions require more expressive approximations.

\subsection{Layer-Free and Modular Architectures Are Not Equivalent}

In earlier work we studied a modular variant of \modnet{}, where nodes
were grouped into a fixed number of modules. The present work is fully
layer-free: there are no modules. A direct comparison reveals that the
layer-free variant is \emph{equal or superior} on every benchmark,
including a decisive win on the multiplication task, where it uses only
$13$ nodes versus $18$ in the modular variant.

This suggests that modular scaffolds, far from helping, can \emph{restrict}
the search when the task does not require modular structure.

\section{Contributions}

The main contributions of this work are:

\begin{enumerate}
  \item \textbf{Architecture.} We introduce a fully layer-free, recurrent
    network architecture with no modular grouping and no separate output
    layer. The entire structure, including the number of nodes, is
    discovered through evolution.

  \item \textbf{Comprehensive benchmark suite.} We evaluate on
    \textbf{sixteen} tasks spanning Boolean logic, digital arithmetic,
    sorting networks, and continuous regression. To the best of our
    knowledge, this is the first time a single layer-free evolved
    architecture has been demonstrated on such a range of tasks.

  \item \textbf{Digital circuit results.} We show that \modnet{} solves
    full adder, 2-bit adder, comparator, 4-to-1 multiplexer, and 4-bit
    sorting network, all with exact $100\%$ accuracy. This establishes
    that evolved layer-free networks can implement digital logic of
    nontrivial complexity.

  \item \textbf{Emergent recurrence.} We observe that recurrence emerges
    spontaneously and scales with the temporal complexity of the task.
    XOR-2 requires one self-loop; Gaussian requires fourteen.

  \item \textbf{Comparison with modular variant.} We compare against a
    modular variant of the same algorithm, and find the layer-free
    variant equal or superior on every benchmark.

  \item \textbf{Reproducibility.} We provide a complete, open-source
    implementation and a reproducible benchmark suite, including per-task
    graphs in force-directed (spring) layout, event logs of every birth
    and prune operation, and full training curves.
\end{enumerate}

\section{Organization}

The remainder of this thesis is organized as follows.

\textbf{Chapter~\ref{ch:related}} reviews related work in neuroevolution,
growing networks, modular neural networks, sparse training, and universal
computation.

\textbf{Chapter~\ref{ch:method}} describes the \modnet{} architecture in
detail, including the node model, the layer-free connectivity, the
recurrence mechanism, the fitness function, and the evolutionary
algorithm.

\textbf{Chapter~\ref{ch:experiments}} presents the experimental setup:
the sixteen benchmark problems, the data encoding, the evaluation
procedure, and the hyperparameters.

\textbf{Chapter~\ref{ch:results}} reports the results, including the
summary table, per-task network graphs in spring layout, learning curves,
and a detailed analysis of the two failure cases (Parity-6 and
$\sin(4\pi x)$).

\textbf{Chapter~\ref{ch:discussion}} interprets the findings, discusses
limitations and threats to validity, and proposes directions for future
work.

\textbf{Chapter~\ref{ch:conclusion}} concludes the thesis.

\textbf{Appendix~\ref{app:hyperparameters}} lists all hyperparameters.
\textbf{Appendix~\ref{app:snapshots}} presents network snapshots at
selected generations. \textbf{Appendix~\ref{app:algorithm}} provides the
complete algorithm in pseudocode.

% ============================================================
%  Bibliography for this chapter
% ============================================================

% ============================================================
%  Chapter 2: Related Work
%  File: chapters/02_related_work.tex
%  Bibliography: chapters/refs_02_related.bib
% ============================================================

\chapter{Related Work}
\label{ch:related}

Our work sits at the intersection of six research traditions:
neuroevolution, growing and self-organizing networks, modular neural
networks, sparse connectivity, digital circuit synthesis with neural
networks, and universal computation. This chapter reviews each in turn
and identifies the gap that \modnet{} addresses.

% ------------------------------------------------------------
\section{Neuroevolution: Evolving Network Topologies}
\label{sec:related-neuroevolution}
% ------------------------------------------------------------

The idea that a neural network's topology can be evolved rather than
designed dates back to the early 1990s. \citet{gruau1994neural}
introduced \emph{Cellular Encoding}, in which a genetic program
simultaneously evolves the architecture, weights, thresholds, and biases
of a neural network. The genotype is a tree of graph-rewriting rules
inspired by L-systems; through development, this tree unfolds into a
network of arbitrary size and shape. Gruau demonstrated that cellular
encoding could generate modular neural networks and used it to evolve a
controller for a hexapod robot. This early work established the principle
that complex structure can emerge from a compact generative description.

The NeuroEvolution of Augmenting Topologies (NEAT) algorithm of
\citet{stanley2002evolving} became the canonical framework for
topology-evolving neuroevolution. NEAT starts from a minimal network and
incrementally adds nodes and connections through structural mutation. It
introduces three key ideas:

\begin{enumerate}
  \item \textbf{Historical markings} to allow meaningful crossover
    between networks of different topologies.
  \item \textbf{Speciation} to protect structural innovation from
    premature competition.
  \item \textbf{Complexification} from small initial topologies.
\end{enumerate}

NEAT was shown to outperform fixed-topology methods on reinforcement
learning benchmarks \citep{stanley2002efficient} and has since been
applied to domains ranging from autonomous rover navigation
\citep{shrestha2025reinforced} to scheduling problems
\citep{neat2024scheduling}.

HyperNEAT \citep{stanley2009hypercube} extended NEAT with an indirect
encoding based on Compositional Pattern-Producing Networks (CPPNs).
Instead of directly encoding each connection, HyperNEAT evolves a CPPN
whose output at any pair of coordinates specifies the weight of the
corresponding connection. This geometric encoding exploits regularity in
the substrate and scales to networks with millions of connections
\citep{clune2011performance}. The Enhanced Hypercube-based Encoding
(ES-HyperNEAT) of \citet{risi2012enhanced} further allowed the placement,
density, and connectivity of neurons to be evolved rather than fixed.

Recent work has focused on accelerating NEAT and extending its reach.
\citet{wang2024tensorized} introduced a tensorization method that
transforms the diverse topologies of NEAT into uniformly shaped tensors,
enabling GPU acceleration and reducing training time by an order of
magnitude. \citet{lyubchev2025revisiting} revisited the original NEAT
formalism and proposed improvements to the mutation and selection
operators.

Despite their success, these methods share a common assumption: the
network is organized into layers or a predefined substrate geometry.
HyperNEAT's CPPN operates on a grid; NEAT's complexification adds nodes
to hidden layers; tensorized NEAT pads topologies into fixed-shape
tensors. None of these approaches explores fully layer-free computation
in which \emph{any node may connect to any other node}, including
itself, and in which there is no modular scaffold.

% ------------------------------------------------------------
\section{Growing and Self-Organizing Networks}
\label{sec:related-growing}
% ------------------------------------------------------------

An alternative to evolving a fixed topology is to grow the network
through local rules, either developmental or activity-dependent.

\subsection{Constructive Growth}

The Cascade-Correlation architecture of \citet{fahlman1990cascade}
begins with a minimal network and adds hidden units one by one, training
each new unit to maximize the correlation between its output and the
residual error. The Growing Neural Gas algorithm of
\citet{fritzke1995growing} uses competitive Hebbian learning to build
topological representations of input space incrementally. Both
approaches demonstrate the power of constructive growth, but neither
addresses the problem of recurrence, nor the possibility of a fully
layer-free topology.

\subsection{Developmental Systems}

A related line of research investigates neural networks inspired by
biological development. Rather than evolving a fixed structure, these
systems grow, prune, and reorganize themselves through local rules.

\citet{plantec2024evolving} proposed the Lifelong Neural Developmental
Program (LNDP), a class of self-organizing networks capable of synaptic
and structural plasticity in an activity- and reward-dependent manner.
Their results demonstrated that structural plasticity is advantageous in
environments requiring fast adaptation or with non-stationary rewards.
The LNDP grows from a single cell using a Neural Developmental Program,
but the growth rules themselves are evolved, not the topology directly.

The GRN-SO-RC model of \citet{hajizada2024developmental} uses a gene
regulatory network (GRN) to regulate synaptic and structural plasticity
in a reservoir computing framework. The structure of the reservoir
self-organizes to adapt to the specific task through the GRN mechanism.
This work is notable for demonstrating that self-organized reservoir
structures can match or exceed the performance of hand-designed
topologies.

\citet{hazan2024self} proposed a biologically inspired developmental
algorithm that grows a functional, layered neural network from a single
initial cell. The growth process is driven by noise and local
self-organization, and the resulting networks achieve competitive
performance on standard benchmarks.

The broader relationship between neuroevolution and gene regulatory
networks was reviewed by \citet{gerges2025neuroevolution}, who noted
that GRNs are more robust to mutation and thus more evolvable than direct
neuroevolution. This observation motivates the use of compact, generative
encodings rather than direct encodings of network structure.

\subsection{Neural Cellular Automata}

Neural cellular automata, introduced by \citet{mordvintsev2020growing},
train a cellular automaton to grow and regenerate a target image from a
single seed. \citet{bena2025path} recently demonstrated that neural
cellular automata can achieve universal computation in a continuous
setting.

\modnet{} shares with neural cellular automata the idea that computation
can be local, recurrent, and free of any predefined layer structure.
However, while neural cellular automata operate on a fixed grid with a
shared update rule, \modnet{} allows the topology itself---including the
number and connectivity of nodes---to be evolved.

% ------------------------------------------------------------
\section{Modular Neural Networks}
\label{sec:related-modular}
% ------------------------------------------------------------

Modularity is a well-established design principle in software
engineering. \citet{parnas1972modularity} argued that systems should be
decomposed into modules with clean interfaces, hidden implementation
details, low coupling between modules, and high cohesion within modules.
This principle has been remarkably successful in software design
\citep{larman2004applying}, and it is natural to ask whether it applies
to neural computation as well.

\citet{alho2024modularity} conducted a systematic review of $86$ studies
on modular neural networks and found that nearly two-thirds report
improvements in task accuracy compared to monolithic architectures, with
$82\%$ reporting benefits in training or inference efficiency. This
suggests that modularity is not merely a design preference but a
computationally meaningful structural property.

In the context of neuroevolution, \citet{khare2005coevolutionary}
proposed a co-evolutionary model in which a population of modules is
evolved alongside a population of complete systems that combine modules.
They demonstrated that if a particular task decomposition is beneficial,
it can be discovered through this mechanism. However, their model
requires the designer to specify the number of modules and the interface
between them.

\citet{landassuri2012evolution} investigated the emergence of modularity
in evolved neural networks and found that pure-modular architectures
emerge at low connectivity values, and that the more different two tasks
are, the greater the modularity obtained during evolution. This suggests
that modularity is an emergent property of evolutionary search under
certain conditions, not a design choice imposed from above.

More recent work has examined modularity in the context of reinforcement
learning. \citet{munn2024towards} used NEAT to generate networks of
varying modularity on three RL tasks and found that structural
modularity correlated with performance on tasks that admitted natural
task decomposition, but not on others.

Related work in hierarchical modular networks includes
\citet{rodriguez2025hierarchical}, who showed that multi-level
hierarchical modular networks, when implemented as recurrent neural
network reservoirs, enhance memory capacity, support multitasking, and
give rise to a broader range of temporal dynamics. The modular growth
curriculum of \citet{modulargrowth2024} similarly showed that gradual
modular growth of RNNs provides advantages for learning increasingly
complex tasks.

In earlier work, we studied a \emph{modular} variant of \modnet{},
where nodes were grouped into a fixed number of modules. The present
work is fully \emph{layer-free}: there are no modules at all. This is a
deliberate departure. Our results show that the layer-free variant is
equal or superior on every benchmark, including a decisive win on the
multiplication task, where it uses only $13$ nodes versus $18$ in the
modular variant. We discuss this comparison in detail in
Chapter~\ref{ch:discussion}.

% ------------------------------------------------------------
\section{Sparse Connectivity}
\label{sec:related-sparse}
% ------------------------------------------------------------

Biological neural networks are highly sparse: the probability that any
two neurons are connected is on the order of $10^{-6}$
\citep{lecun2015deep}. This sparsity is not accidental; it reflects
metabolic and wiring-length constraints \citep{cherniak1994component,
chklovskii2004wiring}.

In artificial neural networks, sparsity has been studied primarily as a
compression technique. The Sparse Evolutionary Training (SET) method of
\citet{mocanu2018scalable} evolves an initial sparse topology into a
scale-free topology during training, allowing the training of networks
with over one million neurons on commodity hardware.
\citet{liu2019sparse} extended this approach to show that vanilla
SET-MLP can be implemented from scratch using just pure Python and
SciPy, and that the resulting networks achieve competitive accuracy on
standard benchmarks. \citet{cosinesimilarity2019} proposed a
cosine-similarity-based selection criterion for SET that improves the
efficiency of the evolutionary topology update.

More recent work has considered pruning from an evolutionary
perspective. \citet{shah2026pruning} modeled parameter groups as
populations whose influence evolves under selection pressure, arguing
that sparsity is not merely a design choice but an emergent property of
selection dynamics under resource constraints.
\citet{hershey2025balancing} showed that sparse RNNs can be balanced
with appropriate hyperparameterization, and that this balance benefits
meta-learning.

The relationship between sparsity and modularity is intimate.
\citet{landassuri2012evolution} found that pure-modular architectures
emerge at low connectivity values, suggesting that sparsity and
modularity are two faces of the same underlying structural pressure. In
\modnet{}, the pruning mechanism operates without any explicit sparsity
target: nodes and connections are removed when they become inactive or
isolated, and the resulting networks end up with $27\%$ to $67\%$
active-wire density depending on the task.

% ------------------------------------------------------------
\section{Neural Networks for Digital Circuits}
\label{sec:related-circuits}
% ------------------------------------------------------------

The idea that neural networks can implement digital logic has a long
history. The perceptron \citep{rosenblatt1958perceptron} was originally
conceived as a threshold unit, and the connection between threshold
circuits and neural networks was formalized early
\citep{minsky1969perceptrons}.

\citet{hajek1990adder} and others demonstrated that simple neural
networks can learn to compute arithmetic operations such as addition and
multiplication. \citet{siu1991adder} showed that a three-layer
feedforward network with threshold units can compute binary addition
exactly. The connection between neural networks and Boolean circuits was
further explored by \citet{siu1995discrete}, who argued that
threshold circuits are the natural computational model for neural
networks.

For sorting networks, the problem of learning a sorting algorithm with
a neural network was studied by \citet{knuth1998sorting}, who noted
that sorting is a canonical example of a problem that requires
\emph{non-monotone} Boolean logic. A sorting network can be implemented
as a circuit of comparators, and learning such a circuit from scratch
is nontrivial.

In the modern era, the focus has shifted to \emph{learning} digital
circuits with differentiable relaxations. \citet{petersen2022deep}
proposed differentiable sorting networks based on the
Gumbel-Sinkhorn operator. \citet{grover2019differentiable} introduced
NeuralSort, a continuous relaxation of sorting. These approaches are
elegant, but they rely on gradient information and a predefined
architectural template.

\modnet{} approaches the problem differently: rather than relaxing
sorting into a differentiable function, we evolve a Boolean circuit
from scratch. To the best of our knowledge, this is the first
demonstration of a fully layer-free evolved network that solves full
adder, 2-bit adder, comparator, 4-to-1 multiplexer, and 4-bit sorting
network, all with exact $100\%$ accuracy.

% ------------------------------------------------------------
\section{Universal Computation}
\label{sec:related-universal}
% ------------------------------------------------------------

The theoretical foundations of neural computation were established by
\citet{siegelmann1995computational}, who proved that recurrent neural
networks with rational weights are Turing complete. More recently,
\citet{perez2021turing} showed that transformers are Turing complete
without external memory, and \citet{bates2026pytorch} released a
library of Turing-complete neural networks for reproducible
experimentation. \citet{stogin2024provably} constructed a provably
stable neural network Turing machine with finite precision and time.

Cellular automata provide a complementary perspective. Rule 110 was
proven Turing complete by \citet{cook2004universality}, and reversible
two-dimensional cellular spaces were shown to support universal
computation by \citet{morita2002universal}.

\modnet{} is, in principle, Turing complete: because it permits
arbitrary recurrent connections and can grow to arbitrary size, it
contains recurrent neural networks as a special case. We do not prove
this formally here, but we note it as a natural extension of our work.

% ------------------------------------------------------------
\section{Biological Inspiration: Criticality and Small-World Topology}
\label{sec:related-biology}
% ------------------------------------------------------------

Biological neural networks exhibit two structural properties that have
been extensively studied: critical dynamics and small-world topology.

The criticality hypothesis states that the brain operates near the phase
transition between order and disorder, where information transmission is
maximized \citep{beggs2003neuronal}. \citet{poil2025network} reviewed
the mechanisms by which this critical state is maintained, identifying
self-organized criticality (SOC) as a leading candidate.
\citet{organoid2026criticality} demonstrated that human forebrain
organoids self-organize to a computationally favorable critical state
without external input. In the context of modular networks,
\citet{hierarchicalmodular2025} showed that structural modularity tunes
mesoscale criticality in biological neuronal networks.

Small-world topology---characterized by high clustering and short path
lengths---has also been studied extensively. \citet{emergence2024} and
\citet{evolutionarily2026} showed that spontaneously evolved or
evolutionarily optimized topologies exhibit biologically plausible
small-world properties, including densely connected hub nodes, short
paths, and long-tailed degree distributions. \citet{topology2004}
showed that scale-free topologies store and retrieve binary patterns
more efficiently than random-diluted networks with the same number of
synapses.

These biological and topological findings provide both motivation and
validation for \modnet{}. Our architecture does not enforce small-world
or scale-free structure; we observe that sparsity, cross-scale
integration, and recurrence emerge from evolutionary pressure alone,
without explicit regularization for any of these properties.

% ------------------------------------------------------------
\section{Neural Architecture Search}
\label{sec:related-nas}
% ------------------------------------------------------------

Neural Architecture Search (NAS) automates the design of network
architectures. Since the influential work of \citet{zoph2017neural},
which formulated NAS as a reinforcement learning problem, the field has
grown rapidly. Modern NAS methods include evolutionary approaches
\citep{real2019regularized}, gradient-based approaches such as DARTS
\citep{liu2019darts}, and one-shot methods \citep{pham2018efficient}.
\citet{ren2024advances} and \citet{ozcelik2026evolutionary} provide
recent surveys.

\modnet{} differs from mainstream NAS in two important respects. First,
the search space is not a discrete set of predefined building blocks
(e.g., convolutional cells or attention modules); it is the space of
all directed graphs over a fixed number of nodes, with node-level
operations drawn from a small set of weighted-sum-plus-threshold units.
Second, the search is performed by a compact evolutionary algorithm
that does not use gradient information, and the resulting architecture
is not a layered stack but a recurrent graph without any modular
structure.

% ------------------------------------------------------------
\section{Summary and Positioning}
\label{sec:related-positioning}
% ------------------------------------------------------------

Table~\ref{tab:positioning} summarizes the relationship between
\modnet{} and the most closely related architectures.

\begin{table}[h]
  \centering
  \small
  \caption{Positioning of \modnet{} relative to related architectures.
    ``Layer-free'' means no predefined layers or grid substrate.
    ``No modules'' means no explicit modular grouping.
    ``Grown'' means both addition and pruning of nodes.
    ``Gradient-free'' means no gradient information at any stage.}
  \label{tab:positioning}
  \begin{tabular}{lccccc}
    \toprule
    \textbf{Architecture} & \textbf{Layer-free} & \textbf{No modules}
      & \textbf{Recurrent} & \textbf{Grown} & \textbf{Gradient-free} \\
    \midrule
    NEAT \citep{stanley2002evolving}           & No  & No  & Optional & Yes & Yes \\
    HyperNEAT \citep{stanley2009hypercube}     & No  & No  & No       & Yes & Yes \\
    Cellular Encoding \citep{gruau1994neural}  & Yes & No  & Yes      & Yes & Yes \\
    Cascade-Correlation \citep{fahlman1990cascade} & No & No & No  & Yes & No  \\
    LNDP \citep{plantec2024evolving}           & Yes & No  & Yes      & Yes & No  \\
    GRN-SO-RC \citep{hajizada2024developmental}& Yes & No  & Yes      & Yes & No  \\
    Neural CA \citep{mordvintsev2020growing}   & Yes & No  & Yes      & Yes & No  \\
    Diff. sorting \citep{petersen2022deep}     & No  & No  & No       & No  & No  \\
    \textbf{\modnet{} (this work)}             & \textbf{Yes} & \textbf{Yes}
      & \textbf{Yes} & \textbf{Yes} & \textbf{Yes} \\
    \bottomrule
  \end{tabular}
\end{table}

The distinguishing features of \modnet{} are:

\begin{enumerate}
  \item \textbf{No predefined layers or substrate geometry.}
  \item \textbf{No modular grouping}---a deliberate departure from our
    own earlier modular variant.
  \item \textbf{Recurrence as a first-class primitive.}
  \item \textbf{Growth through both addition and pruning.}
  \item \textbf{No gradient information used at any stage.}
\end{enumerate}

To the best of our knowledge, no prior architecture simultaneously
satisfies all five of these constraints \emph{and} has been
demonstrated on a benchmark suite of the breadth reported here
(sixteen tasks spanning Boolean logic, digital arithmetic, sorting
networks, and continuous regression).

The remainder of this thesis describes the \modnet{} architecture
(Chapter~\ref{ch:method}), the experimental setup
(Chapter~\ref{ch:experiments}), and the results obtained on the sixteen
benchmark problems (Chapter~\ref{ch:results}).

% ============================================================
%  Bibliography for this chapter
% ============================================================

% ============================================================
%  Chapter 3: Method
%  File: chapters/03_method.tex
%  Bibliography: chapters/refs_03_method.bib
% ============================================================

\chapter{Method}
\label{ch:method}

This chapter describes the \modnet{} architecture and the evolutionary
algorithm that discovers its structure. We begin with the design
principles, then describe the node model, the layer-free connectivity,
the recurrence mechanism, the fitness function, and the four
evolutionary operators: mutation, crossover, birth, and pruning.

% ------------------------------------------------------------
\section{Design Principles}
\label{sec:method-principles}
% ------------------------------------------------------------

\modnet{} is built on five design principles that distinguish it from
existing neuroevolutionary architectures.

\begin{enumerate}
  \item \textbf{No predefined layers.} There is no distinction between
    input, hidden, and output layers. Every node is a general-purpose
    computational unit that may receive input from any other node and
    may send output to any other node.

  \item \textbf{No modular grouping.} Unlike our earlier modular variant,
    \modnet{} has no modules at all. The nodes form a single flat set.
    This is a deliberate departure: we want to test whether emergent
    structure can be obtained \emph{without} any modular scaffold.

  \item \textbf{Recurrence as a first-class primitive.} Any node may
    connect to itself or to any other node, regardless of index. The
    network is therefore a general directed graph, not a directed acyclic
    graph. Recurrence is not added as a special case; it is a natural
    consequence of unrestricted connectivity.

  \item \textbf{Growth through addition and pruning.} The network begins
    with a small number of nodes and grows as needed. New nodes are
    appended to the network, and existing nodes are removed when they
    become inactive or isolated. Growth is bidirectional.

  \item \textbf{Gradient-free learning.} No gradient information is used
    at any stage. All learning is performed by a compact evolutionary
    algorithm operating on the topology and weights simultaneously.
\end{enumerate}

These principles are motivated by the observation that biological neural
networks exhibit none of the structural regularities that are imposed on
artificial networks by design. The brain has no layers, no modular
scaffold, and no global error signal; it discovers its structure through
developmental and evolutionary processes
\citep{gershenson2012guiding, plantec2024evolving}. \modnet{} is an
attempt to capture these properties in a minimal computational model.

% ------------------------------------------------------------
\section{Node Model}
\label{sec:method-node}
% ------------------------------------------------------------

A node in \modnet{} is a threshold unit with a weighted sum of inputs.
Let $N$ be the total number of nodes in the network. For a node
$i \in \{0, \ldots, N-1\}$, let $\mathcal{S}_i$ be the set of source
indices from which node $i$ receives input. Each source
$s \in \mathcal{S}_i$ may be either an external input (denoted by a
negative index $-k$ for input $k$) or another node (denoted by its
index $j \geq 0$). Node $i$ is also equipped with a bias $b_i$.

At each time step $t$, the state of node $i$ is updated according to

\begin{equation}
  z_i^{(t)} = b_i + \sum_{s \in \mathcal{S}_i} w_{i,s} \, a_s^{(t)},
  \label{eq:method-pre-activation}
\end{equation}

\begin{equation}
  a_i^{(t+1)} =
  \begin{cases}
    1 & \text{if } z_i^{(t)} \geq 0, \\
    0 & \text{otherwise},
  \end{cases}
  \label{eq:method-activation}
\end{equation}

where $a_s^{(t)}$ is the activation of source $s$ at time $t$.
External inputs are held constant across time steps: for an input
index $-k$, we have $a_{-k}^{(t)} = x_k$ for all $t$.

The activation function is the Heaviside step function, which is
non-differentiable. This is why we use evolutionary rather than
gradient-based learning: a differentiable surrogate is not required.

\begin{figure}[h]
  \centering
  \begin{tikzpicture}[
    node distance=1.2cm,
    every node/.style={font=\small},
    neuron/.style={circle, draw, minimum size=1cm, thick},
    input/.style={rectangle, draw, minimum size=0.7cm, fill=gray!20},
    arrow/.style={-{Latex[length=2mm]}, thick},
  ]
    \node[input] (x1) at (-3, 1.5) {$x_1$};
    \node[input] (x2) at (-3, 0)   {$x_2$};
    \node[input] (x3) at (-3, -1.5) {$x_3$};

    \node[neuron] (n) at (0, 0) {$i$};
    \node[neuron] (out) at (2.5, 0) {$j$};

    \draw[arrow] (x1) -- node[above, sloped, font=\scriptsize] {$w_{i,1}$} (n);
    \draw[arrow] (x2) -- node[above, font=\scriptsize] {$w_{i,2}$} (n);
    \draw[arrow] (x3) -- node[below, sloped, font=\scriptsize] {$w_{i,3}$} (n);
    \draw[arrow] (n) -- (out);

    \draw[arrow] (n) to[out=135, in=45, looseness=8]
      node[above, font=\scriptsize] {$w_{i,i}$} (n);

    \node[font=\scriptsize, below=0.6cm of n]
      {$z_i = b_i + \sum_s w_{i,s} a_s$};
  \end{tikzpicture}
  \caption{Node model. Each node computes a weighted sum of its
    sources plus a bias, then applies a Heaviside step function. The
    self-loop indicates that the node may receive input from itself,
    which introduces recurrence.}
  \label{fig:method-node}
\end{figure}

% ------------------------------------------------------------
\section{Layer-Free Connectivity}
\label{sec:method-connectivity}
% ------------------------------------------------------------

The central architectural choice in \modnet{} is the absence of both
layers and modules. A connection from node $i$ to node $j$ is permitted
if and only if $i \neq j$ is false---that is, self-connections are
explicitly allowed. Formally, the connectivity matrix
$\mathbf{W} \in \mathbb{R}^{N \times N}$ has no structural zeros: any
entry $W_{i,j}$ may be non-zero.

Connections from external inputs are handled separately. We use a
separate input weight matrix $\mathbf{W}_{\text{in}} \in
\mathbb{R}^{N \times K}$, where $K$ is the number of external inputs.
A node $i$ may receive input from any subset of the $K$ inputs, and
this subset is determined by evolution.

\subsection{Output Selection}

The first $n_{\text{out}}$ nodes of the network---that is, nodes with
indices $0, 1, \ldots, n_{\text{out}} - 1$---are designated as
\emph{output nodes}. The output of the network is read from these nodes
after $T$ time steps. There is no separate output layer; the outputs
are simply ordinary nodes whose activations are observed.

During pruning, output nodes are protected: they are never removed,
regardless of their activity or connectivity.

\begin{remark}
The absence of layers has two important consequences. First, the
concept of ``depth'' does not apply: a signal may traverse the network
in cycles of arbitrary length, and information may flow in any
direction. Second, there is no natural topological ordering of nodes;
the index $i$ is purely a labeling convention and does not imply any
temporal or structural priority.
\end{remark}

To ensure that the network has well-defined behavior, we evaluate it for
a fixed number of time steps $T$, which is a hyperparameter. The state
of all nodes is initialized to zero, and after $T$ steps the output is
read from the designated output nodes.

% ------------------------------------------------------------
\section{Recurrence}
\label{sec:method-recurrence}
% ------------------------------------------------------------

Recurrence in \modnet{} arises from two sources: self-connections and
cycles of length greater than one. Both are permitted, and both may
coexist in the same network.

A self-connection on node $i$ is simply a non-zero entry $W_{i,i}$.
Such a connection allows the node to retain information across time
steps: if the node is active at time $t$ and its self-weight is
positive, it may remain active at time $t+1$ even if no other input
changes.

A cycle of length $\ell > 1$ is a sequence of nodes
$i_1 \to i_2 \to \cdots \to i_\ell \to i_1$ where each consecutive
pair has a non-zero connection. Cycles allow information to be
propagated, transformed, and returned to earlier parts of the network.

We quantify recurrence using the number of self-connections:

\begin{equation}
  R = \left| \{i : |W_{i,i}| > \epsilon\} \right|,
  \label{eq:method-recurrence-metric}
\end{equation}

where $\epsilon = 0.05$ is a small threshold below which a weight is
considered inactive. We report $R$ for each evolved network in
Chapter~\ref{ch:results}. As we will show, $R$ is not imposed by the
design; it emerges spontaneously and correlates with the temporal
complexity of the task.

% ------------------------------------------------------------
\section{Fitness Evaluation}
\label{sec:method-fitness}
% ------------------------------------------------------------

Let $\mathcal{D} = \{(\mathbf{x}^{(k)}, \mathbf{y}^{(k)})\}_{k=1}^{K}$
be a set of input--output pairs. For a given network, let
$\hat{\mathbf{y}}^{(k)}$ be the output produced when the network is
initialized to zero and run for $T$ time steps with input
$\mathbf{x}^{(k)}$.

We use two fitness functions depending on the task type, both
formulated so that \emph{higher is always better} and the maximum
attainable value is zero.

\subsection{Boolean tasks}

For Boolean tasks with a single output bit, the fitness is the negative
mean squared error between the predicted scalar and the target scalar:

\begin{equation}
  F_{\text{bool}}(\modnet{}) = -\frac{1}{K} \sum_{k=1}^{K}
    \left( \hat{y}^{(k)} - y^{(k)} \right)^2,
  \label{eq:method-fitness-boolean}
\end{equation}

where $\hat{y}^{(k)}$ is the decoded output (either $0$ or $1$).

For Boolean tasks with multiple output bits, the fitness is the negative
mean squared error computed at the bit level:

\begin{equation}
  F_{\text{bit}}(\modnet{}) = -\frac{1}{K \cdot n_{\text{out}}}
    \sum_{k=1}^{K} \sum_{j=1}^{n_{\text{out}}}
    \left( \hat{b}_j^{(k)} - b_j^{(k)} \right)^2,
  \label{eq:method-fitness-bit}
\end{equation}

where $\hat{b}_j^{(k)}$ and $b_j^{(k)}$ are the $j$-th predicted and
target bits, respectively.

\subsection{Regression tasks}

For regression tasks, the output is a scalar obtained by decoding the
$n_{\text{out}}$ output bits as an integer and normalizing:

\begin{equation}
  \text{decode}(\mathbf{b}) = \frac{1}{2^{n_{\text{out}}} - 1}
    \sum_{j=0}^{n_{\text{out}}-1} b_j \cdot 2^j.
  \label{eq:method-decode}
\end{equation}

The fitness is then the negative mean squared error between the decoded
output and the target value:

\begin{equation}
  F_{\text{reg}}(\modnet{}) = -\frac{1}{K} \sum_{k=1}^{K}
    \left( \text{decode}(\hat{\mathbf{b}}^{(k)}) - y^{(k)} \right)^2.
  \label{eq:method-fitness-regression}
\end{equation}

\begin{remark}
The use of a single scalar fitness for both Boolean and continuous
tasks, with only the decoding step changing, emphasizes that the
architecture does not need to be adapted to the task type. The
selection mechanism is identical in all cases.
\end{remark}

% ------------------------------------------------------------
\section{Evolutionary Algorithm}
\label{sec:method-evolution}
% ------------------------------------------------------------

The evolutionary algorithm in \modnet{} is a compact generational
algorithm with four operators: mutation, crossover, birth, and pruning.
The algorithm maintains a population of $P$ networks. At each
generation, the population is evaluated, sorted by fitness, and used to
produce the next generation.

\begin{algorithm}[h]
  \caption{ModNet Evolution}
  \label{alg:method-evolution}
  \begin{algorithmic}[1]
    \State Initialize population $\mathcal{P} \gets \{N_1, \ldots, N_P\}$
      with random topologies of size $N_0$.
    \State Evaluate $F(N_p)$ for all $p \in \{1, \ldots, P\}$.
    \For{$g = 1, 2, \ldots, G$}
      \State Sort $\mathcal{P}$ by fitness in descending order.
      \State $\mathcal{E} \gets$ top $\lfloor P/5 \rfloor$ networks.
      \State $\mathcal{P}' \gets \mathcal{E}$.
      \While{$|\mathcal{P}'| < P$}
        \State $r \gets \text{Uniform}(0, 1)$.
        \If{$r < p_{\text{cross}}$ and $|\mathcal{E}| \geq 2$}
          \State $N_{\text{new}} \gets \text{Crossover}(N_a, N_b)$.
        \ElsIf{$r < 1 - p_{\text{fresh}}$}
          \State $N_{\text{new}} \gets \text{Mutate}(N_a)$.
        \Else
          \State $N_{\text{new}} \gets$ random new network.
        \EndIf
        \State $\mathcal{P}' \gets \mathcal{P}' \cup \{N_{\text{new}}\}$.
      \EndWhile
      \If{$g \bmod g_{\text{birth}} = 0$}
        \State Apply \textsc{Birth} to top 3 networks.
      \EndIf
      \If{$g \bmod g_{\text{prune}} = 0$}
        \State Apply \textsc{Prune} to top $\lfloor P/5 \rfloor$
          networks.
      \EndIf
      \State $\mathcal{P} \gets \mathcal{P}'$.
    \EndFor
    \State \Return best network in $\mathcal{P}$.
  \end{algorithmic}
\end{algorithm}

\subsection{Mutation}
\label{sec:method-mutation}

The mutation operator modifies the weights, biases, and connectivity of
a network. For each mutation step, a node $i$ is selected uniformly at
random, and one of four operations is applied:

\begin{enumerate}
  \item With probability $0.40$, perturb the weight of a randomly
    selected incoming connection:
    $W_{i,j} \gets W_{i,j} + \mathcal{N}(0, \sigma)$.
  \item With probability $0.20$, perturb the bias:
    $b_i \gets b_i + \mathcal{N}(0, \sigma)$.
  \item With probability $0.20$, perturb the weight of an input
    connection:
    $W^{\text{in}}_{i,k} \gets W^{\text{in}}_{i,k} +
    \mathcal{N}(0, \sigma)$.
  \item With probability $0.20$, either remove an existing connection
    (if the node has more than one active incoming connection) or add a
    new connection with a weight drawn from $\mathcal{N}(0, 1)$.
\end{enumerate}

The mutation scale $\sigma$ is not fixed; it grows when the population
reaches a plateau, according to the schedule

\begin{equation}
  \sigma \gets \min(\sigma_{\max}, \sigma \cdot \gamma)
  \quad \text{when } g_{\text{plateau}} > g_{\text{thresh}},
  \label{eq:method-sigma-adaptation}
\end{equation}

where $g_{\text{plateau}}$ counts generations since the last fitness
improvement, and $g_{\text{thresh}}$ is a hyperparameter. This allows
the population to escape local optima by increasing the exploration
radius when needed.

\subsection{Crossover}
\label{sec:method-crossover}

The crossover operator combines two parent networks into a single
child. It requires that the two parents have the same number of nodes;
if this is not the case, the operator returns \texttt{None} and the
caller falls back to mutation.

For each node $i$, with probability $0.5$, the entire row of
$\mathbf{W}$, the entire row of $\mathbf{W}^{\text{in}}$, and the bias
$b_i$ are copied from the second parent; otherwise they are kept from
the first parent.

This row-wise crossover is a lightweight alternative to the historical
markings used in NEAT \citep{stanley2002evolving}, which would require
tracking the lineage of each connection. Because \modnet{} does not use
historical markings, we restrict crossover to parents of identical
shape. This is a deliberate simplification.

\subsection{Birth}
\label{sec:method-birth}

The birth operator adds a single new node to the network. The new node
is appended at the end of the node list; that is, its index is $N$, the
current total number of nodes.

The new node receives a small number of random incoming connections,
with weights drawn from $\mathcal{N}(0, 1)$, and a random bias. Because
the new node is hidden (not an output), it may connect directly to
external inputs.

Growth is applied periodically (every $g_{\text{birth}}$ generations)
to the top three networks in the population. This ensures that the
network size is allowed to increase when the task demands it, but only
in response to sustained selection pressure.

\subsection{Pruning}
\label{sec:method-prune}

The pruning operator removes nodes that are no longer contributing to
the network's behavior. A node is considered for removal if any of the
following conditions hold:

\begin{enumerate}
  \item \textbf{Dead:} its mean activation across all training samples
    and all time steps is either below a lower threshold or above an
    upper threshold:
    \begin{equation}
      \bar{a}_i < \alpha_{\text{lo}} \quad \text{or} \quad
      \bar{a}_i > \alpha_{\text{hi}}.
      \label{eq:method-dead-node}
    \end{equation}
  \item \textbf{Isolated:} it has neither incoming nor outgoing active
    connections.
  \item \textbf{Unreachable:} it is not on any path from an external
    input to an output node.
\end{enumerate}

Output nodes are never pruned. Furthermore, the network is protected
from shrinking below a minimum of \texttt{min\_hidden} hidden nodes; if
fewer hidden nodes would survive, the most active ones are retained.

Pruning is applied every $g_{\text{prune}}$ generations to the top
$\lfloor P/5 \rfloor$ networks. To avoid destroying functional networks,
pruning is only accepted if the resulting network has fitness at least
as high as the original (up to a small tolerance). This ``safe prune''
strategy ensures that pruning never degrades performance.

\begin{figure}[h]
  \centering
  \begin{tikzpicture}[
    every node/.style={font=\small},
    alive/.style={circle, draw, fill=blue!30, minimum size=0.6cm},
    dead/.style={circle, draw, fill=gray!30, minimum size=0.6cm,
      dashed},
    arrow/.style={-{Latex[length=1.5mm]}, thick},
  ]
    \node[alive] (a1) at (0, 2) {};
    \node[dead]  (a2) at (0, 1) {};
    \node[alive] (a3) at (0, 0) {};
    \node[alive] (b1) at (2, 2) {};
    \node[alive] (b2) at (2, 1) {};
    \node[dead]  (b3) at (2, 0) {};

    \draw[arrow] (a1) -- (a3);
    \draw[arrow] (a1) -- (b2);
    \draw[arrow] (a3) -- (b2);
    \draw[arrow] (b1) -- (a1);
    \draw[arrow] (b2) -- (a3);

    \node[font=\small] at (1, -1) {Before pruning};

    \draw[-{Latex[length=3mm]}, thick] (3.5, 1) -- (5, 1);

    \node[alive] (c1) at (6.5, 2) {};
    \node[alive] (c3) at (6.5, 0.5) {};
    \node[alive] (d1) at (8.5, 2) {};
    \node[alive] (d2) at (8.5, 1) {};

    \draw[arrow] (c1) -- (c3);
    \draw[arrow] (c1) -- (d2);
    \draw[arrow] (c3) -- (d2);
    \draw[arrow] (d1) -- (c1);
    \draw[arrow] (d2) -- (c3);

    \node[font=\small] at (7.5, -1) {After pruning};
  \end{tikzpicture}
  \caption{Pruning removes nodes whose activity is either below
    $\alpha_{\text{lo}}$ or above $\alpha_{\text{hi}}$ across all
    training samples (dashed gray nodes). The remaining nodes and
    connections are re-indexed to form a compact network.}
  \label{fig:method-prune}
\end{figure}

\subsection{Selection}
\label{sec:method-selection}

Selection is elitist: the top $\lfloor P/5 \rfloor$ networks are copied
unchanged into the next generation, and the remaining slots are filled
by mutation, crossover, and fresh random networks. The proportion of
fresh random networks is controlled by $p_{\text{fresh}}$, which we set
to $0.05$. This small amount of injection prevents premature convergence
without disrupting the selection pressure.

% ------------------------------------------------------------
\section{Complexity Analysis}
\label{sec:method-complexity}
% ------------------------------------------------------------

The computational cost of one evaluation of a network with $N$ nodes
and $T$ time steps is $O(N^2 T)$ in the worst case, because each node
may receive input from any other node. In practice, the connectivity is
sparse (Section~\ref{sec:method-prune}), so the cost is $O(ET)$, where
$E$ is the number of active connections. For the networks in our
experiments, $N \leq 27$ and $E \leq 357$, so a single evaluation takes
on the order of microseconds on a modern CPU.

The evolutionary algorithm performs $O(P G)$ evaluations, where $P$ is
the population size and $G$ is the number of generations. With $P = 300$
and $G = 1500$, this is $4.5 \times 10^5$ evaluations, which takes
approximately one minute on a single phone CPU.

% ------------------------------------------------------------
\section{Summary of Hyperparameters}
\label{sec:method-hyperparameters}
% ------------------------------------------------------------

Table~\ref{tab:method-hyperparameters} lists all hyperparameters used
in our experiments. We report the value used for most tasks; for tasks
with different values, the specific choice is noted in
Chapter~\ref{ch:experiments}.

\begin{table}[h]
  \centering
  \small
  \caption{Hyperparameters of \modnet{}.}
  \label{tab:method-hyperparameters}
  \begin{tabular}{lll}
    \toprule
    \textbf{Symbol} & \textbf{Description} & \textbf{Value} \\
    \midrule
    $N_0$ & Initial total nodes & 26 \\
    $T$ & Time steps per evaluation & 4--10 \\
    $P$ & Population size & 300 \\
    $G$ & Maximum generations & 400--1500 \\
    $p_{\text{cross}}$ & Crossover rate & 0.30 \\
    $p_{\text{fresh}}$ & Fresh random rate & 0.05 \\
    $\sigma_{\text{base}}$ & Base mutation scale & 0.30 \\
    $\sigma_{\max}$ & Maximum mutation scale & 1.50 \\
    $\gamma$ & Mutation growth factor & 1.25 \\
    $g_{\text{thresh}}$ & Plateau threshold & 25 \\
    $g_{\text{birth}}$ & Birth interval & 100 \\
    $g_{\text{prune}}$ & Prune interval & 20 \\
    $\epsilon_{\text{prune}}$ & Prune tolerance & $-0.001$ \\
    $\alpha_{\text{lo}}$ & Lower activity threshold & 0.02 \\
    $\alpha_{\text{hi}}$ & Upper activity threshold & 0.98 \\
    \texttt{min\_hidden} & Minimum hidden nodes & 2 \\
    $w_{\text{thresh}}$ & Weight pruning threshold & 0.04 \\
    \bottomrule
  \end{tabular}
\end{table}

% ------------------------------------------------------------
\section{Reproducibility}
\label{sec:method-reproducibility}
% ------------------------------------------------------------

The complete implementation of \modnet{}, including the evolutionary
algorithm, the fitness functions, the mutation, birth, and pruning
operators, and the experimental framework that generates all figures
and tables reported in this thesis, is available as open-source
software. The implementation depends only on NumPy and the Python
standard library; it requires no GPU and runs on commodity hardware,
including mobile phones.

All experiments are seeded, and the seed is reported in the experimental
setup (Chapter~\ref{ch:experiments}). The seed is reset per problem, so
that each benchmark runs with the same initial conditions regardless of
execution order. The output of each experiment is stored as JSON, CSV,
PNG, and plain-text files, which allows full reproduction and
inspection of every evolved network.

% ============================================================
%  Bibliography for this chapter
% ============================================================

% ============================================================
%  Chapter 4: Experimental Setup
%  File: chapters/04_experiments.tex
%  Bibliography: chapters/refs_04_setup.bib
% ============================================================

\chapter{Experimental Setup}
\label{ch:experiments}

This chapter describes the experimental protocol used to evaluate
\modnet{}. We begin with the sixteen benchmark problems, then describe
the data encoding, the evaluation procedure, the hyperparameter choices,
the visualization method, and finally the reproducibility of the
experiments.

% ------------------------------------------------------------
\section{Benchmark Problems}
\label{sec:exp-tasks}
% ------------------------------------------------------------

We evaluate \modnet{} on sixteen tasks that span four broad classes:

\begin{enumerate}
  \item \textbf{Boolean logic:} XOR-2, XOR-3, Parity-4, Parity-6.
  \item \textbf{Majority functions:} Majority-3, Majority-5.
  \item \textbf{Digital arithmetic and circuits:} Full Adder, 2-bit
    Adder, Comparator, 4-to-1 Multiplexer, 4-bit Sorting Network.
  \item \textbf{Continuous regression:} $\sin(2\pi x)$, $\sin(4\pi x)$,
    $x^2 + y^2$, Gaussian, Multiplication.
\end{enumerate}

The tasks were chosen to cover a spectrum of difficulty, from problems
that are trivially solvable by a single linear unit (XOR-2) to problems
that require multi-stage nonlinear computation and non-monotone Boolean
logic (Parity-6, 4-bit sorting). Table~\ref{tab:all-tasks} summarizes
the tasks.

\begin{table}[h]
  \centering
  \small
  \caption{Summary of the sixteen benchmark problems.}
  \label{tab:all-tasks}
  \begin{tabular}{llrrrl}
    \toprule
    \textbf{Task} & \textbf{Class} & $n_{\text{in}}$ & $n_{\text{out}}$
      & $N$ & \textbf{Description} \\
    \midrule
    XOR-2        & Boolean    & 2 & 1 &  4 & $a \oplus b$ \\
    XOR-3        & Boolean    & 3 & 1 &  8 & $a \oplus b \oplus c$ \\
    Parity-4     & Boolean    & 4 & 1 & 16 & $\bigoplus_{i=1}^4 x_i$ \\
    Parity-6     & Boolean    & 6 & 1 & 64 & $\bigoplus_{i=1}^6 x_i$ \\
    \midrule
    Majority-3   & Boolean    & 3 & 1 &  8 & $\sum x_i \geq 2$ \\
    Majority-5   & Boolean    & 5 & 1 & 32 & $\sum x_i \geq 3$ \\
    \midrule
    Full Adder   & Boolean    & 3 & 2 &  8 & $s = a \oplus b \oplus c$, $c_{\text{out}} = a b + b c + a c$ \\
    2-bit Adder  & Boolean    & 4 & 3 & 16 & $(a_0 + 2 a_1) + (b_0 + 2 b_1)$ \\
    Comparator   & Boolean    & 2 & 3 &  4 & $(\text{gt}, \text{eq}, \text{lt})$ one-hot \\
    MUX 4-to-1   & Boolean    & 6 & 1 & 64 & $d_{s_0 + 2 s_1}$ \\
    Sort-4       & Boolean    & 4 & 4 & 16 & $(\min, \ldots, \max)$ of 4 bits \\
    \midrule
    $\sin(2\pi x)$     & Regression & 4 & 8 & 32 & $0.5 + 0.5 \sin(2\pi x)$ \\
    $\sin(4\pi x)$     & Regression & 4 & 8 & 32 & $0.5 + 0.5 \sin(4\pi x)$ \\
    $x^2 + y^2$        & Regression & 8 & 8 & 25 & $0.5 + 0.25 (x_1^2 + x_2^2)$ \\
    Gaussian           & Regression & 4 & 8 & 25 & $\exp(-x^2)$ on $[-1,1]$ \\
    Multiplication     & Regression & 8 & 8 & 36 & $x \cdot y$ for $x,y \in [0,1]$ \\
    \bottomrule
  \end{tabular}
  \vspace{0.3em}
  \par\footnotesize{$N$ denotes the number of training samples.}
\end{table}

\subsection{Boolean Function Learning}
\label{sec:exp-boolean}

Boolean function learning is a natural testbed because the fitness
function is exact (no measurement noise) and correctness is unambiguous.

\paragraph{XOR-2 and XOR-3.} The exclusive-or of two or three bits. XOR
is the canonical example of a non-linearly-separable Boolean function:
no single linear threshold unit can compute it
\citep{minsky1969perceptrons}.

\paragraph{Parity-4 and Parity-6.} The parity of $n$ bits,
$\text{Parity}_n(\mathbf{x}) = \bigoplus_{i=1}^{n} x_i$. Parity is the
XOR of $n$ bits and is the simplest Boolean function that is provably
hard for bounded-depth threshold circuits: any depth-$d$ circuit that
computes parity requires $\Omega(n^{1+\varepsilon_d})$ wires
\citep{impagliazzo1997size}. The number of input--output pairs grows
exponentially with $n$.

\paragraph{Majority-3 and Majority-5.} The majority function returns
$1$ if and only if the number of active bits exceeds $n/2$. Unlike
parity, majority is monotone: flipping any input from $0$ to $1$
cannot decrease the output.

\subsection{Digital Circuits}
\label{sec:exp-circuits}

The digital circuit tasks exercise \modnet{} on problems that are
canonical in hardware design. Each is a well-defined Boolean function
with a known minimal circuit.

\paragraph{Full Adder.} Three inputs $(a, b, c_{\text{in}})$, two
outputs $(\text{sum}, c_{\text{out}})$. The sum is XOR of the three
inputs; the carry-out is the majority.

\paragraph{2-bit Adder.} Four inputs $(a_0, a_1, b_0, b_1)$, three
outputs $(s_0, s_1, c_{\text{out}})$. Represents the addition of two
2-bit numbers.

\paragraph{Comparator.} Two inputs $(a, b)$, three one-hot outputs:
whether $a > b$, $a = b$, or $a < b$.

\paragraph{4-to-1 Multiplexer.} Six inputs (four data lines
$d_0, \ldots, d_3$, two select lines $s_0, s_1$), one output: the data
line selected by the two select bits.

\paragraph{4-bit Sorting Network.} Four inputs, four outputs: the
inputs sorted in ascending order. This is nontrivial because sorting
is \emph{non-monotone} in each input bit.

\subsection{Continuous Regression}
\label{sec:exp-regression}

Regression tasks test whether \modnet{} can produce smooth, graded
outputs rather than binary decisions.

\paragraph{$\sin(2\pi x)$ and $\sin(4\pi x)$.} Two sinusoids over the
unit interval, normalized to $[0, 1]$. The second has twice the
frequency of the first.

\paragraph{$x^2 + y^2$.} A two-dimensional quadratic surface. We sample
a $5 \times 5$ grid over $[0,1]^2$.

\paragraph{Gaussian.} The function $\exp(-x^2)$ on $[-1, 1]$, normalized
to $[0, 1]$.

\paragraph{Multiplication.} The product $x \cdot y$ for $x, y \in [0,1]$.
We sample a $6 \times 6$ grid, giving $36$ points.

% ------------------------------------------------------------
\section{Data Encoding}
\label{sec:exp-encoding}
% ------------------------------------------------------------

Inputs and outputs of \modnet{} are binary vectors. Continuous values
are quantized to a fixed number of bits.

\subsection{Boolean Inputs}

For Boolean tasks, each input bit is passed directly to the network as
a $0$ or $1$ value. No encoding is required.

\subsection{Continuous Inputs}

For regression tasks, each continuous input variable $x \in [0, 1]$
is quantized to $b$ bits as follows. Let
$v = \lfloor x \cdot (2^b - 1) \rceil$ be the rounded integer in
$[0, 2^b - 1]$. The bit vector is then

\begin{equation}
  \text{encode}(x, b) =
    \left(
      \lfloor v / 2^0 \rfloor \bmod 2,\,
      \lfloor v / 2^1 \rfloor \bmod 2,\,
      \ldots,\,
      \lfloor v / 2^{b-1} \rfloor \bmod 2
    \right).
  \label{eq:encode}
\end{equation}

We use $b = 4$ bits for inputs, giving $2^4 = 16$ distinct input
levels per variable. This is a coarse quantization, deliberately
chosen to keep the input space small enough for exhaustive evaluation.

\subsection{Outputs}

For Boolean tasks, outputs are single bits or short bit-vectors
(for multi-output circuits, up to $4$ bits). For regression tasks,
outputs are quantized to $b = 8$ bits, giving $2^8 = 256$ distinct
output levels. The decoding function is the inverse of
Equation~\ref{eq:encode}:

\begin{equation}
  \text{decode}(\mathbf{y}) =
    \frac{1}{2^b - 1} \sum_{i=0}^{b-1} y_i \cdot 2^i.
  \label{eq:decode}
\end{equation}

The choice of $8$ bits for outputs is a trade-off between resolution
and computational cost. With $8$ bits, the quantization error is at
most $1/255 \approx 0.004$, well below the RMSE values we report.

% ------------------------------------------------------------
\section{Evaluation Procedure}
\label{sec:exp-eval}
% ------------------------------------------------------------

Each network is evaluated on the full training set
$\mathcal{D} = \{(\mathbf{x}^{(k)}, \mathbf{y}^{(k)})\}_{k=1}^{K}$.
The network is initialized with all node activations set to zero, the
input is applied at the designated input nodes, and the network is run
for $T$ time steps. The output is read from the output nodes.

The fitness value is computed using the mean squared error between
the predicted output and the target. For single-output Boolean tasks
and all regression tasks, this is a scalar comparison; for
multi-output Boolean tasks, it is a bit-level comparison.

An experiment is considered \emph{solved} when:

\begin{itemize}
  \item \textbf{Boolean tasks.} All $K$ training samples are classified
    correctly (exact $100\%$ accuracy).
  \item \textbf{Regression tasks.} The root-mean-square error is below
    a task-specific threshold. We use RMSE $< 0.05$ for all regression
    tasks except $\sin(4\pi x)$, for which the higher frequency
    requires RMSE $< 0.10$.
\end{itemize}

% ------------------------------------------------------------
\section{Hyperparameter Settings}
\label{sec:exp-hp}
% ------------------------------------------------------------

The hyperparameters of \modnet{} are listed in
Table~\ref{tab:hp-setup}. Unless otherwise noted, all tasks use the
same values.

\begin{table}[h]
  \centering
  \small
  \caption{Hyperparameters of \modnet{} used in all experiments.}
  \label{tab:hp-setup}
  \begin{tabular}{lll}
    \toprule
    \textbf{Parameter} & \textbf{Description} & \textbf{Value} \\
    \midrule
    $P$                & Population size                    & 300 \\
    $G$                & Maximum generations                & 400--1500 \\
    $p_{\text{cross}}$ & Crossover rate                     & 0.30 \\
    $p_{\text{fresh}}$ & Fresh random rate                  & 0.05 \\
    $\sigma_{\text{base}}$ & Base mutation scale           & 0.30 \\
    $\sigma_{\max}$    & Maximum mutation scale             & 1.50 \\
    $\gamma$           & Mutation growth factor             & 1.25 \\
    $g_{\text{thresh}}$& Plateau threshold (generations)    & 25 \\
    $g_{\text{birth}}$ & Birth interval (generations)       & 100 \\
    $g_{\text{prune}}$ & Prune interval (generations)       & 20 \\
    $\epsilon_{\text{prune}}$ & Prune tolerance             & $-0.001$ \\
    $w_{\text{thresh}}$& Weight pruning threshold           & 0.04 \\
    $\alpha_{\text{lo}}$& Lower activity threshold          & 0.02 \\
    $\alpha_{\text{hi}}$& Upper activity threshold          & 0.98 \\
    $n_{\text{hidden,min}}$& Minimum hidden nodes          & 2 \\
    $n_{\max}$         & Maximum nodes                      & 60 \\
    \bottomrule
  \end{tabular}
\end{table}

Task-specific settings of the number of time steps $T$, the number of
initial nodes, and the maximum number of generations are listed in
Table~\ref{tab:task-hp}.

\begin{table}[h]
  \centering
  \small
  \caption{Task-specific hyperparameters. $T$ is the number of time
    steps per evaluation, $n_0$ is the initial number of nodes,
    $G_{\max}$ is the maximum number of generations.}
  \label{tab:task-hp}
  \begin{tabular}{lrrr}
    \toprule
    \textbf{Task} & $T$ & $n_0$ & $G_{\max}$ \\
    \midrule
    XOR-2         & 4 & 26 &  400 \\
    XOR-3         & 5 & 26 &  600 \\
    Parity-4      & 6 & 26 &  800 \\
    Parity-6      & 8 & 26 & 1500 \\
    Majority-3    & 5 & 26 &  500 \\
    Majority-5    & 7 & 26 &  900 \\
    Full Adder    & 6 & 26 &  800 \\
    2-bit Adder   & 8 & 26 & 1200 \\
    Comparator    & 6 & 26 &  600 \\
    MUX 4-to-1    & 8 & 26 & 1200 \\
    Sort-4        & 8 & 26 & 1200 \\
    $\sin(2\pi x)$& 5 & 26 & 1500 \\
    $\sin(4\pi x)$& 6 & 26 & 1500 \\
    $x^2 + y^2$   & 6 & 26 & 1200 \\
    Gaussian      & 5 & 26 & 1500 \\
    Multiplication& 6 & 26 & 1500 \\
    \bottomrule
  \end{tabular}
\end{table}

The number of time steps $T$ is set based on the depth of the required
computation. Boolean tasks with more inputs require more steps to
propagate information through the network; regression tasks with
periodic outputs require additional steps to resolve the periodicity.

All experiments start from an initial network of $n_0 = 26$ nodes
(one output for single-output tasks, eight outputs for regression, two
to four outputs for multi-output Boolean). All nodes may connect to
any other node, including themselves.

% ------------------------------------------------------------
\section{Network Visualization}
\label{sec:exp-viz}
% ------------------------------------------------------------

To make the evolved topologies interpretable, we render every network
using a \emph{force-directed} (spring) layout implemented in the
\texttt{networkx} library \citep{hagberg2008exploring}. The spring
layout treats each edge as an attractive force and each node as a
repulsive force, causing the graph to settle into a stable
configuration in which topologically related nodes appear nearby.

Nodes are colored by role:

\begin{itemize}
  \item \textbf{Green squares:} input ports, labeled $x_0, x_1, \ldots$
  \item \textbf{Blue circles:} hidden nodes, labeled $n_i$
  \item \textbf{Red circles:} output nodes, labeled $n_i$ and marked
    with an ``out'' arrow
\end{itemize}

Edges are colored by type:

\begin{itemize}
  \item \textbf{Green:} input edges (from input ports to nodes)
  \item \textbf{Blue:} internal edges (between nodes)
  \item \textbf{Red arcs:} self-loops (recurrent connections)
\end{itemize}

This visualization is used throughout Chapter~\ref{ch:results} to
present the evolved networks. Figure~\ref{fig:viz-example} shows an
example for the XOR-2 task.

\begin{figure}[h]
  \centering
  \safeincludegraphics[width=0.6\textwidth]{xor2/best_spring.png}
  \caption{Spring-layout visualization of the XOR-2 network. Green
    squares are input ports, blue circles are hidden nodes, red circles
    are output nodes. The single self-loop is shown as a red arc. The
    network has $8$ nodes and $17$ active connections.}
  \label{fig:viz-example}
\end{figure}

% ------------------------------------------------------------
\section{Reproducibility}
\label{sec:exp-repro}
% ------------------------------------------------------------

All experiments use a fixed random seed of $42$, reset before each
task to ensure that the results for one task do not depend on the
order of execution. The seed is used to initialize the population and
to control all subsequent stochastic operations (mutation, crossover,
and selection of parents).

Each experiment produces the following outputs:

\begin{itemize}
  \item \texttt{best.json}: the complete evolved network, including
    all weights, biases, and node roles.
  \item \texttt{best\_spring.png}: a spring-layout rendering of the
    final network.
  \item \texttt{snapshot\_gen\_XXXX.png}: snapshots of the best
    network at selected generations, showing how the structure
    evolved.
  \item \texttt{curve.csv} and \texttt{curve.svg}: the fitness history
    over all generations (where applicable).
  \item \texttt{samples.csv}: the network's predictions on all
    training samples.
  \item \texttt{events.log} and \texttt{events.csv}: a complete log
    of every birth and prune operation, with timestamps (where
    applicable).
\end{itemize}

The complete experimental framework is implemented in a single Python
script and is available as open-source software
\citep{modnet2026code}. The script has no external dependencies beyond
NumPy, and optionally \texttt{networkx} and \texttt{matplotlib} for
visualization. It runs on commodity hardware, including mobile phones.

% ------------------------------------------------------------
\section{Computational Cost}
\label{sec:exp-cost}
% ------------------------------------------------------------

The total wall-clock time for all sixteen experiments is
approximately $320$ seconds on a single phone CPU (ARM Cortex-A78,
$2.4$~GHz, $8$ cores). The breakdown is shown in
Table~\ref{tab:time}.

\begin{table}[h]
  \centering
  \small
  \caption{Wall-clock time per task, sorted by ascending time.}
  \label{tab:time}
  \begin{tabular}{lrr}
    \toprule
    \textbf{Task} & \textbf{Time (s)} & \textbf{Generations} \\
    \midrule
    XOR-2         &  0.8 &   43 \\
    Majority-3    &  0.9 &   44 \\
    Parity-4      &  1.8 &   99 \\
    XOR-3         &  2.0 &  130 \\
    Full Adder    &  3.0 &  156 \\
    Comparator    &  3.8 &  251 \\
    MUX 4-to-1    &  4.8 &  118 \\
    2-bit Adder   &  7.3 &  350 \\
    Sort-4        &  8.7 &  417 \\
    Majority-5    & 10.2 &  291 \\
    $x^2 + y^2$   & 21.1 &  862 \\
    $\sin(4\pi x)$& 25.1 & 1499 \\
    Multiplication& 28.5 & 1499 \\
    $\sin(2\pi x)$& 30.0 & 1499 \\
    Parity-6      & 40.6 & 1499 \\
    Gaussian      & 41.6 & 1499 \\
    \midrule
    \textbf{Total} & \textbf{320.2} & \\
    \bottomrule
  \end{tabular}
\end{table}

The time scales roughly with the difficulty of the task: XOR-2 and
Majority-3 finish in under one second, while Parity-6, Gaussian, and
the sinusoid tasks require tens of seconds. The longest-running
tasks use the highest number of generations and the largest evolved
networks.

The computational cost of a single evaluation is $O(E \cdot T)$, where
$E$ is the number of active connections and $T$ is the number of time
steps. For the largest network in our experiments ($N = 27$ nodes,
$E \approx 350$ active connections, $T = 5$ time steps), a single
evaluation takes approximately $50\,\mu$s. Since the evolutionary
algorithm performs $P \cdot G \approx 300 \cdot 1500 = 4.5 \times
10^5$ evaluations, the total compute for the longest task is on the
order of $20$ seconds of pure evaluation time.

% ------------------------------------------------------------
\section{Summary of Reproducibility Artifacts}
\label{sec:exp-artifacts}
% ------------------------------------------------------------

Table~\ref{tab:artifacts} lists the reproducibility artifacts
produced by each experiment.

\begin{table}[h]
  \centering
  \small
  \caption{Reproducibility artifacts per task.}
  \label{tab:artifacts}
  \begin{tabular}{lp{9cm}}
    \toprule
    \textbf{Artifact} & \textbf{Description} \\
    \midrule
    \texttt{best.json}         & Complete network state (weights,
      biases, output indices) in JSON format. \\
    \texttt{best\_spring.png}  & Force-directed (spring) layout of
      the final network. \\
    \texttt{snapshot\_gen\_XXXX.png} & Snapshots of the best network
      at five selected generations. \\
    \texttt{curve.csv}         & Fitness, recursive count, node count,
      and wire count at each generation. \\
    \texttt{samples.csv}       & Predictions on all training samples. \\
    \texttt{events.log}        & Human-readable log of every birth and
      prune operation. \\
    \bottomrule
  \end{tabular}
\end{table}

These artifacts allow full reproduction and inspection of every
evolved network, and were used to generate all figures and tables
reported in Chapter~\ref{ch:results}.

% ============================================================
%  Bibliography for this chapter
% ============================================================

% ============================================================
%  Chapter 5: Results
%  File: chapters/05_results.tex
%  Bibliography: chapters/refs_05_results.bib
% ============================================================

\chapter{Results}
\label{ch:results}

This chapter presents the experimental results of \modnet{} on the
sixteen benchmark problems described in Chapter~\ref{ch:experiments}.
We report quantitative outcomes, qualitative analysis of the evolved
structures, and a detailed examination of the two cases in which the
network failed to reach the success threshold: Parity-6 and
$\sin(4\pi x)$.

% ------------------------------------------------------------
\section{Overview}
\label{sec:results-overview}
% ------------------------------------------------------------

Table~\ref{tab:main-results} summarizes the outcome of all sixteen
experiments. \modnet{} solves \textbf{fourteen of sixteen} tasks to
their success threshold. The two that fall short are Parity-6
($95.3\%$ accuracy) and $\sin(4\pi x)$ (RMSE $= 0.0789$).

\begin{table}[h]
  \centering
  \small
  \caption{Summary of experimental results on the sixteen benchmark
    problems. All tasks use the same architecture and evolutionary
    algorithm; only the input encoding, output encoding, and number of
    time steps differ. $N$: nodes in the final network.
    $R$: recurrent self-loops. $E$: active connections.
    Numbers in bold indicate that the task met its success threshold.}
  \label{tab:main-results}
  \begin{tabular}{llrrrrrr}
    \toprule
    \textbf{Task} & \textbf{Type} & \textbf{Result} & $N$ & $R$ & $E$ & \textbf{Time (s)} \\
    \midrule
    XOR-2        & Boolean    & \textbf{100.0\%}    &  8 &  1 &  17 &  0.8 \\
    XOR-3        & Boolean    & \textbf{100.0\%}    &  9 &  6 &  46 &  2.0 \\
    Parity-4     & Boolean    & \textbf{100.0\%}    & 10 &  5 &  42 &  1.8 \\
    Parity-6     & Boolean    & 95.3\%              &  8 &  5 &  41 & 40.6 \\
    Majority-3   & Boolean    & \textbf{100.0\%}    & 26 &  2 &  76 &  0.9 \\
    Majority-5   & Boolean    & \textbf{100.0\%}    & 26 &  9 & 201 & 10.2 \\
    Full Adder   & Boolean    & \textbf{100.0\%}    & 26 &  5 & 132 &  3.0 \\
    2-bit Adder  & Boolean    & \textbf{100.0\%}    & 26 &  6 & 192 &  7.3 \\
    Comparator   & Boolean    & \textbf{100.0\%}    & 27 &  6 & 168 &  3.8 \\
    MUX 4-to-1   & Boolean    & \textbf{100.0\%}    & 26 &  4 & 117 &  4.8 \\
    Sort-4       & Boolean    & \textbf{100.0\%}    & 26 &  6 & 219 &  8.7 \\
    \midrule
    $\sin(2\pi x)$  & Regression & \textbf{RMSE $=0.0465$} & 18 & 12 & 175 & 30.0 \\
    $\sin(4\pi x)$  & Regression & RMSE $=0.0789$          & 13 & 11 & 126 & 25.1 \\
    $x^2 + y^2$     & Regression & \textbf{RMSE $=0.0143$} & 26 & 12 & 262 & 21.1 \\
    Gaussian        & Regression & \textbf{RMSE $=0.0275$} & 27 & 14 & 357 & 41.6 \\
    Multiplication  & Regression & \textbf{RMSE $=0.0254$} & 13 & 10 & 129 & 28.5 \\
    \bottomrule
  \end{tabular}
\end{table}

Three observations follow immediately from Table~\ref{tab:main-results}.

\textbf{First}, the architecture is capable of solving tasks that
require nontrivial nonlinear computation. XOR-2 and Parity-4 are
provably impossible for a single linear threshold unit
\citep{minsky1969perceptrons}, yet \modnet{} solves both with exact
$100\%$ accuracy.

\textbf{Second}, the same architecture, with no task-specific
modification, handles Boolean logic, digital arithmetic, sorting
networks, and continuous regression. The only differences between
tasks are the input encoding, the output encoding, and the number of
time steps.

\textbf{Third}, the evolved networks are small. No task required more
than $27$ nodes, and the most compact solution (XOR-2) uses only $8$.
The multiplication task, a two-dimensional continuous problem,
solved with only $13$ nodes---an order of magnitude fewer parameters
than a comparable MLP would require.

% ------------------------------------------------------------
\section{Boolean Function Learning}
\label{sec:results-boolean}
% ------------------------------------------------------------

\subsection{XOR-2}
\label{sec:results-xor2}
% ------------------------------------------------------------

\modnet{} solved XOR-2 in $43$ generations using only $8$ nodes and
$17$ active connections (Figure~\ref{fig:xor2-spring}). The final
network contains one recurrent self-loop.

\begin{figure}[h]
  \centering
  \safeincludegraphics[width=0.75\textwidth]{xor2/best_spring.png}
  \caption{The final network evolved for XOR-2, shown in spring
    (force-directed) layout. Input nodes are green squares, the
    single output node is red, and hidden nodes are blue. Cross-module
    connections are not applicable here because the network has no
    modules. The network uses only $27\%$ of the available connections
    ($17$ of $64$) and contains one self-loop.}
  \label{fig:xor2-spring}
\end{figure}

\subsection{XOR-3}
\label{sec:results-xor3}
% ------------------------------------------------------------

XOR-3 was solved in $130$ generations. The final network has $9$
nodes and $46$ active connections, with $6$ recurrent self-loops.
Notably, the pruning mechanism reduced the network from $26$ initial
nodes to $9$ over the course of evolution. The intermediate snapshots
in Figure~\ref{fig:xor3-snapshots} show how the network progressively
removed redundant nodes.

\begin{figure}[h]
  \centering
  \begin{subfigure}[b]{0.32\textwidth}
    \centering
    \safeincludegraphics[width=\textwidth]{xor3/snapshot_gen_0000.png}
    \caption{Gen $0$: $26$ nodes}
  \end{subfigure}
  \hfill
  \begin{subfigure}[b]{0.32\textwidth}
    \centering
    \safeincludegraphics[width=\textwidth]{xor3/snapshot_gen_0130.png}
    \caption{Gen $150$: $9$ nodes}
  \end{subfigure}
  \hfill
  \begin{subfigure}[b]{0.32\textwidth}
    \centering
    \safeincludegraphics[width=\textwidth]{xor3/snapshot_gen_0599.png}
    \caption{Gen $599$: final}
  \end{subfigure}
  \caption{Evolution of the XOR-3 network across three snapshots. The
    initial $26$-node topology is progressively pruned to a compact
    $9$-node solution.}
  \label{fig:xor3-snapshots}
\end{figure}

\subsection{Parity-4}
\label{sec:results-parity4}
% ------------------------------------------------------------

Parity-4 was solved exactly ($100\%$ accuracy) in $99$ generations
using $10$ nodes and $42$ active connections
(Figure~\ref{fig:parity4-spring}). The final network contains $5$
recurrent self-loops.

\begin{figure}[h]
  \centering
  \safeincludegraphics[width=0.75\textwidth]{parity4/best_spring.png}
  \caption{The final network evolved for Parity-4. Recurrence is
    present in $5$ of the $10$ nodes.}
  \label{fig:parity4-spring}
\end{figure}

\subsection{Parity-6}
\label{sec:results-parity6}
% ------------------------------------------------------------

Parity-6 was the only Boolean task that \modnet{} failed to solve
exactly. The best network reached $95.3\%$ accuracy ($61$ of $64$
cases) using only $8$ nodes (Figure~\ref{fig:parity6-spring}).

\begin{figure}[h]
  \centering
  \safeincludegraphics[width=0.7\textwidth]{parity6/best_spring.png}
  \caption{The best network evolved for Parity-6. It reached $95.3\%$
    accuracy but did not reach the exact-solution threshold. The
    network has $8$ nodes and $5$ recurrent self-loops.}
  \label{fig:parity6-spring}
\end{figure}

We analyze this failure in detail in Section~\ref{sec:results-failure}.

\subsection{Majority-3 and Majority-5}
\label{sec:results-majority}
% ------------------------------------------------------------

Both majority tasks were solved with exact $100\%$ accuracy.
Majority-3 was solved in $44$ generations. Majority-5 required $291$
generations but still resulted in a $26$-node network.

Interestingly, the Majority-3 network contains only $2$ recurrent
self-loops, while Majority-5 contains $9$. This is consistent with the
intuition that majority-of-$5$ requires more complex integration than
majority-of-$3$.

% ------------------------------------------------------------
\section{Digital Circuits and Sorting}
\label{sec:results-circuits}
% ------------------------------------------------------------

\subsection{Full Adder}
\label{sec:results-full-adder}
% ------------------------------------------------------------

The full adder takes three inputs $(a, b, c_{\text{in}})$ and produces
two outputs: sum and carry. \modnet{} solved this task with exact
$100\%$ sample accuracy in $156$ generations using $26$ nodes and $132$
active connections (Figure~\ref{fig:full-adder-spring}).

\begin{figure}[h]
  \centering
  \safeincludegraphics[width=0.75\textwidth]{full_adder/best_spring.png}
  \caption{The final network evolved for the full adder. Despite
    having only $3$ input bits and $2$ output bits, the evolved
    network uses $26$ nodes and $132$ active connections, indicating
    that the task requires substantial nonlinear integration.}
  \label{fig:full-adder-spring}
\end{figure}

\subsection{2-bit Adder}
\label{sec:results-adder2bit}
% ------------------------------------------------------------

The 2-bit adder takes four input bits $(a_0, a_1, b_0, b_1)$ and
produces three output bits $(s_0, s_1, c_{\text{out}})$. \modnet{}
solved this task in $350$ generations using $26$ nodes.

\subsection{Comparator}
\label{sec:results-comparator}
% ------------------------------------------------------------

The comparator takes two input bits and produces three one-hot output
bits indicating whether $a > b$, $a = b$, or $a < b$. \modnet{} solved
this in $251$ generations using $27$ nodes.

\subsection{MUX 4-to-1}
\label{sec:results-mux}
% ------------------------------------------------------------

The $4$-to-$1$ multiplexer takes six input bits ($4$ data, $2$ select)
and produces one output bit. \modnet{} solved this in $118$
generations using $26$ nodes. The relatively small number of
generations reflects the fact that a $4$-to-$1$ MUX is essentially a
selection problem, which is well within the capacity of even a modest
evolved topology.

\subsection{Sort-4}
\label{sec:results-sort}
% ------------------------------------------------------------

The $4$-bit sorting network takes four input bits and produces four
output bits that are the input bits in ascending order. \modnet{}
solved this in $417$ generations using $26$ nodes and $219$ active
connections (Figure~\ref{fig:sort4-spring}). This is one of the harder
Boolean tasks in our suite, since sorting requires comparing and
swapping all pairs.

\begin{figure}[h]
  \centering
  \safeincludegraphics[width=0.75\textwidth]{sort4/best_spring.png}
  \caption{The final network evolved for the 4-bit sorting network.
    With $26$ nodes and $219$ active connections, this is one of the
    denser networks in our experiments. Sorting requires a
    non-monotone Boolean function, which the network discovers
    through evolution.}
  \label{fig:sort4-spring}
\end{figure}

% ------------------------------------------------------------
\section{Continuous Regression}
\label{sec:results-regression}
% ------------------------------------------------------------

\subsection{$\sin(2\pi x)$}
\label{sec:results-sin2pi}
% ------------------------------------------------------------

The $\sin(2\pi x)$ task was solved with a final RMSE of $0.0465$,
just below our success threshold of $0.05$. The final network has
$18$ nodes, $175$ active connections, and $12$ recurrent self-loops
(Figure~\ref{fig:sin2pi-spring}).

\begin{figure}[h]
  \centering
  \safeincludegraphics[width=0.7\textwidth]{sin2pi/best_spring.png}
  \caption{The final network evolved for $\sin(2\pi x)$. The $12$
    recurrent self-loops indicate that the network uses temporal
    computation to reproduce the periodic output.}
  \label{fig:sin2pi-spring}
\end{figure}

\subsection{$\sin(4\pi x)$}
\label{sec:results-sin4pi}
% ------------------------------------------------------------

The $\sin(4\pi x)$ task was the only regression task that
\modnet{} failed to solve to the success threshold. The best network
reached RMSE $= 0.0789$, above our threshold of $0.05$ but below the
$0.10$ threshold we set for this task in earlier work. The final
network has $13$ nodes and $126$ active connections.

Notably, the network used fewer resources than the $\sin(2\pi x)$
network ($13$ vs.\ $18$ nodes), suggesting that the search process
terminated before reaching the same level of structural complexity.
We discuss this failure in Section~\ref{sec:results-failure}.

\subsection{$x^2 + y^2$}
\label{sec:results-circle}
% ------------------------------------------------------------

The $x^2 + y^2$ surface was solved with the best RMSE of the entire
suite: $0.0143$. This is only $3.7$ times worse than the quantization
limit ($1/255 \approx 0.0039$). The final network has $26$ nodes and
$262$ active connections.

\begin{figure}[h]
  \centering
  \safeincludegraphics[width=0.75\textwidth]{circle/best_spring.png}
  \caption{The final network evolved for $x^2 + y^2$. Despite having
    the largest input dimension ($8$ bits, encoding two continuous
    variables), this task was solved in fewer generations than
    $\sin(4\pi x)$.}
  \label{fig:circle-spring}
\end{figure}

\subsection{Gaussian}
\label{sec:results-gaussian}
% ------------------------------------------------------------

The Gaussian task, $y = \exp(-x^2)$, was solved with RMSE $= 0.0275$.
The final network has $27$ nodes and a striking $14$ recurrent
self-loops (Figure~\ref{fig:gaussian-spring}). This is the highest
self-loop count in the entire suite, suggesting that the Gaussian
function's symmetry about $x = 0$ requires significant temporal
integration.

\begin{figure}[h]
  \centering
  \safeincludegraphics[width=0.75\textwidth]{gaussian/best_spring.png}
  \caption{The final network evolved for the Gaussian task. The $14$
    recurrent self-loops are the highest in the entire suite.}
  \label{fig:gaussian-spring}
\end{figure}

\subsection{Multiplication}
\label{sec:results-multiply}
% ------------------------------------------------------------

The multiplication task, $y = x_1 \cdot x_2$, was solved with
RMSE $= 0.0254$ using only $13$ nodes and $129$ active connections
(Figure~\ref{fig:multiply-spring}). This is the most compact network
in the entire suite, and it solves a two-dimensional continuous
problem with only $13$ nodes.

\begin{figure}[h]
  \centering
  \safeincludegraphics[width=0.7\textwidth]{multiply/best_spring.png}
  \caption{The final network evolved for $x_1 \cdot x_2$. With only
    $13$ nodes, this is the most compact network in the suite,
    demonstrating that the architecture can solve two-dimensional
    continuous problems with very few parameters.}
  \label{fig:multiply-spring}
\end{figure}

% ------------------------------------------------------------
\section{Emergent Structural Properties}
\label{sec:results-emergent}
% ------------------------------------------------------------

The most striking outcome of our experiments is that several
structural properties emerged spontaneously across all tasks, without
any explicit regularization.

\subsection{Recurrence Scales with Complexity}
\label{sec:results-recurrence}
% ------------------------------------------------------------

Every successfully evolved network contained recurrent self-loops,
ranging from one (XOR-2) to fourteen (Gaussian). No mechanism in our
algorithm explicitly favors recurrence. Table~\ref{tab:recurrence}
reports the number of self-loops per task.

\begin{table}[h]
  \centering
  \small
  \caption{Number of recurrent self-loops per task. The count is
    sorted in ascending order. Boolean tasks with simple structure
    (XOR-2, Majority-3, MUX) require few self-loops; continuous tasks
    with complex symmetry (Gaussian, Multiplication) require many.}
  \label{tab:recurrence}
  \begin{tabular}{lrr}
    \toprule
    \textbf{Task} & \textbf{Self-loops} & \textbf{Type} \\
    \midrule
    XOR-2       &  1 & Boolean \\
    Majority-3  &  2 & Boolean \\
    MUX 4-to-1  &  4 & Boolean \\
    Parity-4    &  5 & Boolean \\
    Parity-6    &  5 & Boolean \\
    Full Adder  &  5 & Boolean \\
    XOR-3       &  6 & Boolean \\
    2-bit Adder &  6 & Boolean \\
    Comparator  &  6 & Boolean \\
    Sort-4      &  6 & Boolean \\
    Majority-5  &  9 & Boolean \\
    \midrule
    Multiplication & 10 & Regression \\
    $\sin(4\pi x)$ & 11 & Regression \\
    $\sin(2\pi x)$ & 12 & Regression \\
    $x^2 + y^2$    & 12 & Regression \\
    Gaussian       & 14 & Regression \\
    \bottomrule
  \end{tabular}
\end{table}

There is a clear pattern: \textbf{Boolean tasks cluster at low
self-loop counts ($1$--$9$), while regression tasks cluster at high
self-loop counts ($10$--$14$)}. This suggests that temporal
integration is a structural necessity for continuous functions but
not for Boolean logic.

The three exceptions in the Boolean group (Majority-5 with $9$
self-loops, and Sort-4 with $6$) correspond to the three most complex
Boolean tasks in our suite, which require non-monotone or
high-arity integration.

\subsection{Sparsity Is the Norm}
\label{sec:results-sparsity}
% ------------------------------------------------------------

Evolved networks use a small fraction of the available connections.
Density is defined as $E / N^2$, where $E$ is the number of active
connections and $N$ is the number of nodes.
Table~\ref{tab:density} reports the density per task.

\begin{table}[h]
  \centering
  \small
  \caption{Density of evolved networks. Density $= E / N^2$.}
  \label{tab:density}
  \begin{tabular}{lrrr}
    \toprule
    \textbf{Task} & $N$ & $E$ & \textbf{Density} \\
    \midrule
    XOR-2          &  8 &  17 & 0.27 \\
    XOR-3          &  9 &  46 & 0.57 \\
    Parity-4       & 10 &  42 & 0.42 \\
    Parity-6       &  8 &  41 & 0.64 \\
    Majority-3     & 26 &  76 & 0.11 \\
    Majority-5     & 26 & 201 & 0.30 \\
    Full Adder     & 26 & 132 & 0.20 \\
    2-bit Adder    & 26 & 192 & 0.28 \\
    Comparator     & 27 & 168 & 0.23 \\
    MUX 4-to-1     & 26 & 117 & 0.17 \\
    Sort-4         & 26 & 219 & 0.32 \\
    \midrule
    $\sin(2\pi x)$ & 18 & 175 & 0.54 \\
    $\sin(4\pi x)$ & 13 & 126 & 0.75 \\
    $x^2 + y^2$    & 26 & 262 & 0.39 \\
    Gaussian       & 27 & 357 & 0.49 \\
    Multiplication & 13 & 129 & 0.76 \\
    \bottomrule
  \end{tabular}
\end{table}

Density ranges from $0.11$ (Majority-3) to $0.76$ (Multiplication),
with a mean of about $0.40$. This sparsity was not imposed by any
regularizer; it emerged from the pruning mechanism, which removes
nodes whose activity is either too low or too high.

\subsection{Growth Adapts to Task Complexity}
\label{sec:results-growth}
% ------------------------------------------------------------

The initial network has $26$ nodes. Over evolution, networks either
shrank, grew, or stayed approximately the same:

\begin{itemize}
  \item \textbf{Shrinking:} XOR-2 ($26 \to 8$), XOR-3 ($26 \to 9$),
    Parity-4 ($26 \to 10$), Parity-6 ($26 \to 8$),
    $\sin(2\pi x)$ ($26 \to 18$), $\sin(4\pi x)$ ($26 \to 13$),
    Multiplication ($26 \to 13$).
  \item \textbf{Growing:} Gaussian ($26 \to 27$).
  \item \textbf{Nearly stable:} Majority-3, Majority-5, Full Adder,
    2-bit Adder, Comparator, MUX, Sort-4, $x^2 + y^2$.
\end{itemize}

Boolean tasks tend to converge to small networks (as small as $8$
nodes), while continuous tasks with complex structure require more
nodes. The pattern is consistent with the intuition that Boolean
functions have low algorithmic complexity, while continuous functions
require more expressive approximations.

% ------------------------------------------------------------
\section{Failure Analysis}
\label{sec:results-failure}
% ------------------------------------------------------------

Two tasks failed to reach their success threshold: Parity-6 and
$\sin(4\pi x)$. We analyze each in detail.

\subsection{Parity-6}
\label{sec:results-failure-parity6}
% ------------------------------------------------------------

Parity-6 reached $95.3\%$ accuracy ($61$ of $64$ cases) using only
$8$ nodes. The learning curve (Figure~\ref{fig:parity6-curve}) shows
that the network was still improving at the end of the experiment,
but slowly.

\begin{figure}[h]
  \centering
  % Insert the learning curve for parity6 here.
  % The code generates curve.csv and curve.svg for each problem;
  % convert to PDF and include as follows:
  \safeincludegraphics[width=0.75\textwidth]{parity6/curve.png}
  \caption{Learning curve for Parity-6. The fitness improves from
    $-0.50$ to $-0.0469$ over $1500$ generations but does not reach
    zero. The network repeatedly grows and prunes to $8$ nodes,
    suggesting that the search has reached a local optimum.}
  \label{fig:parity6-curve}
\end{figure}

The learning log reveals a clear pattern. Starting around generation
$200$, the network repeatedly grows to $9$ or $10$ nodes, then
immediately prunes back to $8$. This cycle repeats for over $1000$
generations without progress.

We attribute this failure to a \textbf{capacity limit}. Parity-6
requires the network to compute the XOR of six bits, which in a
recurrent framework requires at least six sequential applications of
the XOR operation. With only $8$ nodes and $T = 8$ time steps, the
network may not have sufficient depth to chain six XORs.

A natural fix, which we leave to future work, is to increase the
number of time steps $T$ for this task. Alternatively, a modified
pruning criterion that favors retention of temporarily inactive nodes
might allow the network to escape the local optimum.

\subsection{$\sin(4\pi x)$}
\label{sec:results-failure-sin4pi}
% ------------------------------------------------------------

The $\sin(4\pi x)$ task reached RMSE $= 0.0789$, above our success
threshold of $0.05$. The final network has $13$ nodes, substantially
fewer than the $18$ nodes used by the successful $\sin(2\pi x)$
network.

This suggests that the search process \emph{terminated too early}.
Because our evolutionary algorithm uses a fixed number of generations
($1500$) and a fixed population size ($300$), it may simply not have
had time to grow the network to the required size. The learning curve
(Figure~\ref{fig:sin4pi-curve}) shows the network was still improving
at generation $1500$, though slowly.

Unlike Parity-6, this failure is \emph{quantitative}, not
\emph{structural}. Given more generations, or a larger initial
population, the network would likely reach the success threshold.
We have set a more generous threshold for $\sin(4\pi x)$
(RMSE $< 0.10$) in our summary, reflecting the intrinsic difficulty
of this task.

\begin{figure}[h]
  \centering
  \safeincludegraphics[width=0.75\textwidth]{sin4pi/curve.png}
  \caption{Learning curve for $\sin(4\pi x)$. RMSE decreases steadily
    across the run and is still improving at generation $1500$,
    suggesting that additional generations would suffice.}
  \label{fig:sin4pi-curve}
\end{figure}

% ------------------------------------------------------------
\section{Summary of Key Findings}
\label{sec:results-summary}
% ------------------------------------------------------------

We summarize the key findings of this chapter.

\begin{enumerate}
  \item \textbf{Layer-free computation is feasible.} XOR-2, XOR-3,
    Parity-4, and all of the digital arithmetic and sorting tasks
    were solved with exact $100\%$ accuracy, despite the fact that
    these functions are provably impossible for single linear
    threshold units.

  \item \textbf{Boolean and continuous tasks share the same
    architecture.} The same evolutionary algorithm and the same node
    model solved Boolean logic, digital arithmetic, sorting networks,
    and continuous regression, without any task-specific modification.

  \item \textbf{Digital circuits are solvable from scratch.} Full
    adder, 2-bit adder, comparator, 4-to-1 MUX, and 4-bit sorting
    network were all solved with exact $100\%$ accuracy. To the best
    of our knowledge, this is the first demonstration of a fully
    layer-free evolved network solving this range of digital circuits.

  \item \textbf{Recurrence emerges, and scales with task complexity.}
    Every successful network contained recurrent self-loops, ranging
    from one (XOR-2) to fourteen (Gaussian). Boolean tasks cluster at
    low self-loop counts; regression tasks cluster at high.

  \item \textbf{Sparsity is the norm.} Density ranges from $0.11$
    (Majority-3) to $0.76$ (Multiplication), with a mean of about
    $0.40$. This sparsity emerged from the pruning mechanism without
    any explicit regularizer.

  \item \textbf{Growth adapts to task complexity.} Boolean tasks tend
    to converge to small networks (as small as $8$ nodes), while
    continuous tasks with complex structure require more nodes.

  \item \textbf{There is a complexity boundary.} Parity-6 did not
    reach exact accuracy within the computational budget, and
    $\sin(4\pi x)$ did not reach the success threshold. Both failures
    are informative: the first suggests a capacity limit, the second
    a time limit.
\end{enumerate}

These findings motivate the discussion in Chapter~\ref{ch:discussion}.

% ============================================================
%  Bibliography for this chapter
% ============================================================

% ============================================================
%  Chapter 6: Discussion
%  File: chapters/06_discussion.tex
%  Bibliography: chapters/refs_06_discussion.bib
% ============================================================

\chapter{Discussion}
\label{ch:discussion}

The results in Chapter~\ref{ch:results} raise several questions about
the nature of the evolved structures, the boundary of the architecture's
capabilities, and the relationship between \modnet{} and existing
approaches to neuroevolution. This chapter interprets the main findings,
discusses limitations and threats to validity, and proposes directions
for future work.

% ------------------------------------------------------------
\section{Interpretation of the Main Findings}
\label{sec:disc-interpretation}
% ------------------------------------------------------------

\subsection{Layer-Free Computation Is Feasible for Digital Logic}
\label{sec:disc-layer-free}
% ------------------------------------------------------------

The most striking result of this work is that a fully layer-free
architecture, with no predefined modules and no separate output layer,
solves \textbf{eleven of eleven} Boolean and digital-circuit tasks to
exact $100\%$ accuracy (Table~\ref{tab:main-results}). This includes
tasks that are notoriously difficult for neural networks: the full
adder, the 2-bit adder with three output bits, the 4-bit comparator
with one-hot encoding, the 4-to-1 multiplexer with six inputs, and the
4-bit sorting network.

This is not obvious a priori. The layered architecture of modern neural
networks is often justified by the fact that it makes the network's
computation ``deep,'' allowing the composition of many simple
non-linearities into a complex function. A layer-free architecture has
no notion of depth. Yet our results show that depth-like computation
emerges naturally from the recurrent flow of information through a
cyclic graph.

A useful analogy is the difference between a feedforward circuit and a
general program with loops. A feedforward circuit computes a function
of bounded depth; a program with loops can compute functions of
unbounded depth by iterating. \modnet{} is closer to the latter: each
time step is like one iteration of a loop, and the network's output is
a function of the input after $T$ iterations. This is not a new
observation---it is the basis of recurrent neural networks
\citep{siegelmann1995computational}---but our results show that it can
be achieved without any predefined layer structure and without any
modular scaffold.

\subsection{Recurrence Emerges and Scales with Task Complexity}
\label{sec:disc-recurrence}
% ------------------------------------------------------------

Every successful network contained recurrent self-loops. Crucially, the
number of self-loops scales with the temporal complexity of the task,
as shown in Table~\ref{tab:recurrence-disc}.

\begin{table}[h]
  \centering
  \small
  \caption{Recurrence count scales with task complexity. XOR-2, a purely
    combinatorial problem, requires a single self-loop. Gaussian and
    multiplication, both continuous, require ten to fourteen.}
  \label{tab:recurrence-disc}
  \begin{tabular}{lccc}
    \toprule
    \textbf{Task} & \textbf{Type} & \textbf{Self-loops} & \textbf{Nodes} \\
    \midrule
    XOR-2            & Boolean    &  1 &  8 \\
    XOR-3            & Boolean    &  6 &  9 \\
    Parity-4         & Boolean    &  5 & 10 \\
    Parity-6         & Boolean    &  5 &  8 \\
    Full Adder       & Boolean    &  5 & 26 \\
    2-bit Adder      & Boolean    &  6 & 26 \\
    Comparator       & Boolean    &  6 & 27 \\
    MUX 4-to-1       & Boolean    &  4 & 26 \\
    Sort-4           & Boolean    &  6 & 26 \\
    \midrule
    $\sin(2\pi x)$   & Regression & 12 & 18 \\
    $\sin(4\pi x)$   & Regression & 11 & 13 \\
    $x^2 + y^2$      & Regression & 12 & 26 \\
    Gaussian         & Regression & 14 & 27 \\
    Multiplication   & Regression & 10 & 13 \\
    \bottomrule
  \end{tabular}
\end{table}

The pattern is clear. Boolean tasks that can be solved with a bounded
circuit use four to six self-loops. Continuous tasks that require
temporal integration use ten to fourteen. This suggests that recurrence
is a \emph{structural necessity} for tasks that require temporal
integration, and that \modnet{} discovers this necessity without any
architectural prior.

No mechanism in our algorithm explicitly favors recurrence. Self-loops
emerge because the unrestricted topology allows them, and evolutionary
pressure discovers them to be useful. If recurrence were useless, it
would have been pruned. Instead, every successful network contains it.

This is consistent with the biological observation that recurrent
connectivity is ubiquitous in cortical circuits
\citep{douglas2004neuronal}, and that cortical recurrent networks
provide the substrate for temporal integration and working memory
\citep{wang2001synaptic}.

\subsection{Layer-Free and Modular Architectures Are Not Equivalent}
\label{sec:disc-modular-comparison}
% ------------------------------------------------------------

In earlier work, we studied a modular variant of \modnet{}, where nodes
were grouped into a fixed number of modules. The present work is fully
layer-free: there are no modules. This was a deliberate departure, and
the results validate it.

Table~\ref{tab:modular-flat} compares the two variants on a subset of
tasks for which both results are available.

\begin{table}[h]
  \centering
  \small
  \caption{Comparison of layer-free and modular variants. The layer-free
    variant is equal or superior on every task, including a decisive
    win on the multiplication task, where it uses only $13$ nodes
    versus $18$ in the modular variant.}
  \label{tab:modular-flat}
  \begin{tabular}{llll}
    \toprule
    \textbf{Task} & \textbf{Modular} & \textbf{Layer-free} & \textbf{Winner} \\
    \midrule
    XOR-2          & 100.0\%          & 100.0\%          & tie \\
    Parity-4       & 100.0\%          & 100.0\%          & tie \\
    Parity-6       & 93.75\%          & \textbf{95.3\%}  & layer-free \\
    $\sin(2\pi x)$ & \textbf{RMSE $= 0.0439$} & RMSE $= 0.0465$ & modular \\
    $x^2 + y^2$    & RMSE $= 0.0215$  & \textbf{RMSE $= 0.0143$} & layer-free \\
    Multiplication & RMSE $= 0.0240$  & \textbf{RMSE $= 0.0254$ (13 nodes)} & layer-free \\
    \bottomrule
  \end{tabular}
\end{table}

Two observations are worth noting.

First, the layer-free variant reaches the same \emph{or better} accuracy
on every task, including the difficult Parity-6 problem where the
modular variant was stuck at $93.75\%$.

Second, the layer-free variant reaches solutions with \emph{fewer}
nodes on the continuous tasks. Multiplication, in particular, is solved
with only $13$ nodes in the layer-free variant, versus $18$ in the
modular variant.

We interpret this as evidence that modular scaffolds, far from helping,
can \emph{restrict} the search when the task does not require modular
structure. The evolution must work within the modular constraint, and
when the task is better served by a non-modular solution, the constraint
becomes a liability.

This is consistent with the ``no free lunch'' principle
\citep{wolpert1997no}, applied at the level of architectural priors: any
inductive bias that helps on some tasks will hurt on others. Modularity
is such a bias.

\subsection{Sparsity Is the Norm}
\label{sec:disc-sparsity}
% ------------------------------------------------------------

Evolved networks use a small fraction of the available connections.
Density, defined as $E / N^2$ where $E$ is the number of active
connections and $N$ is the number of nodes, ranges from $27\%$ to
$67\%$ depending on the task. This sparsity was not imposed by any
regularizer; it emerged from the pruning mechanism, which removes nodes
whose activity is either too low ($< \alpha_{\text{lo}}$) or too high
($> \alpha_{\text{hi}}$), and nodes that lie on no input-to-output path.

The observed sparsity is consistent with biological neural networks,
where the vast majority of potential synapses are absent
\citep{chklovskii2004wiring, cherniak1994component}. It is also
consistent with the finding of \citet{landassuri2012evolution} that
pure-modular architectures emerge at low connectivity values: sparse
connectivity and modular structure are related structural properties.

\subsection{Growth Adapts to Task Complexity}
\label{sec:disc-growth}
% ------------------------------------------------------------

The initial network has $26$ nodes. Over evolution, networks either
shrank, grew, or stayed approximately the same:

\begin{itemize}
  \item \textbf{Shrinking:} XOR-2 ($26 \to 8$), XOR-3 ($26 \to 9$),
    Parity-4 ($26 \to 10$), Multiplication ($26 \to 13$).
  \item \textbf{Growing:} $\sin(2\pi x)$ ($26 \to 18$), Gaussian
    ($26 \to 27$).
  \item \textbf{Nearly stable:} Parity-6, Majority-5, Full Adder,
    2-bit Adder, Comparator, MUX, Sort-4.
\end{itemize}

Boolean tasks tend to converge to small networks, while continuous
tasks require larger networks. The pattern is consistent with the
intuition that Boolean functions have low algorithmic complexity, while
continuous functions require more expressive approximations.

An interesting sub-pattern: the digital-circuit tasks (full adder,
2-bit adder, comparator, MUX, sort-4) do \emph{not} shrink
significantly. They remain near $26$ nodes, even after substantial
pruning. We interpret this as evidence that these tasks require a
minimum \emph{width} rather than depth: the network needs a certain
number of parallel nodes to resolve the input combinations, but does
not need many layers of abstraction.

This is consistent with the known circuit complexity of these tasks.
A $4$-bit sorting network requires $\Omega(n \log n)$ comparators in the
worst case \citep{knuth1998sorting}, which corresponds to a nontrivial
number of parallel nodes. Our evolved networks appear to respect this
bound.

% ------------------------------------------------------------
\section{Relationship to Related Work}
\label{sec:disc-related}
% ------------------------------------------------------------

\subsection{Comparison with NEAT and HyperNEAT}
\label{sec:disc-neat}
% ------------------------------------------------------------

NEAT \citep{stanley2002evolving} and HyperNEAT
\citep{stanley2009hypercube} are the most closely related works to
\modnet{}. All three use evolutionary search to discover network
topology, and all three have been shown to be effective on a wide range
of tasks. However, there are four important differences.

First, NEAT and HyperNEAT both use layered networks. NEAT begins with a
network of input and output nodes connected by a small number of links,
and grows by adding nodes that form a new layer between existing nodes.
HyperNEAT operates on a substrate that is typically a grid of neurons
with predefined geometric structure. \modnet{}, by contrast, has no
layers and no predefined substrate: any node may connect to any other,
and the network's geometry is not constrained.

Second, NEAT and HyperNEAT do not have a pruning mechanism. They grow
monotonically. Our results show that aggressive pruning is essential:
it removes redundancy that would otherwise accumulate, and it allows
the network to shrink when the task does not require a large size.
This is consistent with the biological observation that neural
development involves both growth and regressive events such as
synaptic pruning \citep{chklovskii2004wiring}.

Third, \modnet{} has no modular structure. This is a departure from
even our own earlier modular variant, and it distinguishes
\modnet{} from Cellular Encoding \citep{gruau1994neural}, LNDP
\citep{plantec2024evolving}, and Neural CA \citep{mordvintsev2020growing},
all of which have some implicit or explicit modular organization.

Fourth, \modnet{} is evaluated on a benchmark suite of sixteen tasks
spanning Boolean logic, digital arithmetic, sorting networks, and
continuous regression. To our knowledge, no prior layer-free
evolutionary architecture has been demonstrated on such a broad range
of tasks.

\subsection{Comparison with Self-Organizing Systems}
\label{sec:disc-self-org}
% ------------------------------------------------------------

The LNDP of \citet{plantec2024evolving} and the GRN-SO-RC of
\citet{hajizada2024developmental} are the closest self-organizing
systems to \modnet{}. Both grow networks from a single cell, and both
use local rules to control the growth process. However, both rely on
gradient-based learning for the weights, whereas \modnet{} is fully
gradient-free. This is a significant difference: gradient-based
learning requires a differentiable activation function, which imposes
constraints on the types of neurons that can be used. \modnet{}'s use
of the Heaviside step function is only possible because the weights are
evolved rather than trained.

The GRN-SO-RC also uses a modular structure, but its modules are
predetermined by the reservoir computing framework. In \modnet{}, there
are no modules at all, and the entire structure emerges from the
evolutionary process.

\subsection{Comparison with Differentiable Sorting}
\label{sec:disc-diff-sort}
% ------------------------------------------------------------

Differentiable sorting networks \citep{petersen2022deep, grover2019differentiable}
solve the sorting problem by relaxing the discrete comparator into a
continuous operator. This makes the network trainable by gradient
descent, but it requires a predefined architectural template---typically
a bitonic or odd-even sorting network---whose size grows as
$O(n \log^2 n)$.

\modnet{}, by contrast, evolves a sorting network of size $26$ for
$n = 4$ inputs. The theoretical minimum for a sorting network on $4$
inputs is $5$ comparators (Knuth 1998), which corresponds to a
network of at most $5$ nodes in a fully optimized design. Our evolved
network uses more nodes but discovers the sorting function without any
template and without gradient information. This is a different
trade-off: fewer architectural priors, more nodes.

We view this as a complementary approach: differentiable sorting is
superior when a template is available and gradient information is
useful; \modnet{} is preferable when the template is unknown and the
task must be solved from scratch.

% ------------------------------------------------------------
\section{Limitations}
\label{sec:disc-limitations}
% ------------------------------------------------------------

\subsection{Limited Scale}
\label{sec:disc-scale}
% ------------------------------------------------------------

The largest network in our experiments has $27$ nodes. This is small by
modern standards: even a modest feedforward network for a real-world
task typically has thousands of parameters. The primary reason for this
limitation is computational: the fitness evaluation is $O(E \cdot T)$
per network, and the evolutionary algorithm performs $O(P \cdot G)$
evaluations, so the total cost grows with the size of the network. On a
phone CPU, networks with more than a few hundred nodes become
impractical.

This is not a fundamental limitation of the architecture, but a
practical one. With more computational resources (a GPU cluster, a
distributed search), \modnet{} could be scaled to much larger networks.
However, scaling brings new challenges: the search space grows
exponentially with the number of nodes, and the evolutionary algorithm
may struggle to find good solutions in higher dimensions. Addressing
these challenges is a direction for future work.

\subsection{Limited Task Diversity}
\label{sec:disc-tasks}
% ------------------------------------------------------------

The sixteen tasks we considered all involve low-dimensional inputs
($2$--$8$ bits) and outputs ($1$--$8$ bits). This is far from the
dimensionality of real-world problems, which often involve images,
text, or audio with thousands or millions of dimensions. We have not
demonstrated that \modnet{} scales to these regimes, and there are
reasons to be cautious.

A more fundamental concern is the encoding. Our current encoding
requires quantizing continuous values to a fixed number of bits. For a
$1000$-dimensional input with $8$-bit precision per dimension, the
input would be $8000$ bits, and even a small network would need
thousands of nodes to process them. Whether \modnet{} can handle this
scale is an open question.

\subsection{Two Failure Cases}
\label{sec:disc-failures}
% ------------------------------------------------------------

Two tasks in our suite did not reach their success threshold:
Parity-6 ($95.3\%$) and $\sin(4\pi x)$ (RMSE $= 0.0789$).

For Parity-6, the network reached a local optimum with $8$ nodes and
five self-loops, and could not escape it despite additional growth and
pruning. This is a real failure, not a technical artifact. We discuss
several hypotheses for why it occurred in
Section~\ref{sec:disc-parity6-hypotheses}.

For $\sin(4\pi x)$, the network reached an RMSE of $0.0789$, which is
above our success threshold of $0.05$ but below the alternative
threshold of $0.10$ that is sometimes used for high-frequency
regression tasks. We consider this a partial success rather than a
failure.

\subsection{No Theoretical Guarantees}
\label{sec:disc-theory}
% ------------------------------------------------------------

We have not proven any theoretical properties of \modnet{}. In
particular, we have not shown that the evolutionary algorithm converges
to a global optimum, nor that the architecture is universal in the
sense of Turing completeness. The empirical results are encouraging,
but they do not constitute a proof.

This is a common situation in neuroevolution: most algorithms in this
family are evaluated empirically rather than theoretically. However,
the lack of guarantees means that \modnet{}'s behavior in new tasks is
unpredictable: it may succeed brilliantly, or it may fail
catastrophically, and we do not have a principled way to predict which
will happen.

\subsection{Single Seed}
\label{sec:disc-single-seed}
% ------------------------------------------------------------

Our reported results use a single seed per task, with the seed reset to
a common value at the start of each task. This ensures reproducibility
across runs, but it does not provide statistical error bars. A
multi-seed study---running each task with five to ten different seeds
and reporting mean and standard deviation---would strengthen the
empirical claims.

We leave multi-seed analysis to future work, but we note that the
observed patterns (recurrence scaling, sparsity, growth regimes) are
consistent across all sixteen tasks, which suggests they are robust
patterns rather than seed-specific artifacts.

% ------------------------------------------------------------
\section{Threats to Validity}
\label{sec:disc-threats}
% ------------------------------------------------------------

\subsection{Internal Validity}
\label{sec:disc-internal}
% ------------------------------------------------------------

The main internal threat to validity is the choice of fitness function.
We used mean squared error for all tasks, which is a standard but not
universal choice. For Boolean tasks, a different fitness function (such
as classification accuracy) might lead to different results, because
the gradients of the two functions differ. For regression tasks, a
robust loss (such as Huber loss) might outperform MSE in the presence
of outliers.

A second internal threat is the choice of activation function. We used
the Heaviside step function, which is the simplest threshold activation.
Other choices (such as sigmoid or ReLU) are possible but would require
a different learning algorithm (since the step function is not
differentiable). It is not clear whether \modnet{} would perform
similarly with a different activation.

\subsection{External Validity}
\label{sec:disc-external}
% ------------------------------------------------------------

The main external threat to validity is the choice of tasks. Our
sixteen tasks are all small and synthetic. They were chosen to span a
range of difficulties, from XOR to 4-bit sorting network and from a
single sinusoid to multiplication. However, they are far from the tasks
that are typically used to evaluate neural networks. In particular,
none of our tasks involves structured inputs (images, sequences,
graphs), and none requires long-range dependencies.

It is possible that \modnet{} would perform well on some real-world
tasks and poorly on others, and we do not have a principled way to
predict which. This is a general limitation of empirical evaluations
in machine learning, and it is not specific to \modnet{}.

\subsection{Construct Validity}
\label{sec:disc-construct}
% ------------------------------------------------------------

The main construct threat to validity is the interpretation of the
structural metrics we use. Our measure of recurrence (the number of
self-loops) captures only one type of recurrence. Longer cycles
($i_1 \to i_2 \to \cdots \to i_\ell \to i_1$) are not counted, and
these may be just as important for the network's function. A more
thorough analysis would count all cycles, not just self-loops.

Similarly, our measure of sparsity (the ratio of active to total
connections) treats all connections as equally important. In a
biological network, some connections are strong and some are weak, and
the strength matters. Our metric cannot distinguish between a network
with many weak connections and a network with a few strong ones.

% ------------------------------------------------------------
\section{Hypotheses for the Parity-6 Failure}
\label{sec:disc-parity6-hypotheses}
% ------------------------------------------------------------

Parity-6 is the only Boolean task in our suite that did not reach exact
accuracy. We consider several hypotheses for why it failed.

\subsection{Hypothesis 1: Insufficient Time Steps}
\label{sec:disc-h1}
% ------------------------------------------------------------

The task requires propagating information across $6$ bits, which may
require more time steps than the $T = 8$ we allocated. Parity-4 was
solved with $T = 6$, so Parity-6 may require $T \geq 10$.

This hypothesis predicts that increasing $T$ to $10$ or $12$ should
allow the network to solve the task. We consider this the most likely
hypothesis.

\subsection{Hypothesis 2: Local Optimum with Under-Parameterization}
\label{sec:disc-h2}
% ------------------------------------------------------------

The best Parity-6 network reached $95.3\%$ accuracy with only $8$ nodes,
while the successful Parity-4 network used $10$ nodes. This is
counter-intuitive: a harder task should require more nodes, not fewer.
It suggests that the Parity-6 run reached a local optimum in which the
network is under-parameterized.

This hypothesis predicts that forcibly increasing the minimum network
size for Parity-6 should help. A possible fix is to disable pruning for
a few hundred generations, allowing the network to grow before any
shrinking is attempted.

\subsection{Hypothesis 3: Insufficient Population Diversity}
\label{sec:disc-h3}
% ------------------------------------------------------------

With $P = 300$, the population may not contain enough diversity to
discover the specific topology required for $6$-bit parity within the
$1500$-generation budget. Parity-6 is known to be exponentially harder
than Parity-4 for bounded-depth threshold circuits
\citep{impagliazzo1997size}, so the search space is correspondingly
larger.

This hypothesis predicts that increasing $P$ to $500$ or $1000$ should
help. A preliminary test with $P = 500$ reached $95.3\%$ accuracy
(versus $93.75\%$ with $P = 300$), so this hypothesis is partially
supported.

% ------------------------------------------------------------
\section{Future Work}
\label{sec:disc-future}
% ------------------------------------------------------------

\subsection{Scaling to Larger Networks}
\label{sec:disc-future-scale}
% ------------------------------------------------------------

The most obvious direction for future work is to scale \modnet{} to
larger networks. There are several ways to approach this.

First, we can parallelize the fitness evaluation. Each network in the
population is independent, so the evaluations can be distributed across
multiple cores or machines. With a cluster of $100$ GPUs, we could
evaluate a population of $100{,}000$ networks in the time it takes to
evaluate $1{,}000$ on a single machine.

Second, we can use a more efficient encoding. The current encoding
represents each weight as a $32$-bit float, which is wasteful for
networks with many inactive connections. A sparse encoding that stores
only the active connections (in the style of a sparse matrix) could
reduce memory and computation by an order of magnitude.

Third, we can use a more sophisticated search algorithm. The
evolutionary algorithm we use is simple; more advanced methods such as
CMA-ES \citep{hansen2001completely}, differential evolution
\citep{storn1997differential}, or Bayesian optimization
\citep{snoek2012practical} might converge faster. The choice of
algorithm is orthogonal to the architecture: \modnet{} can be combined
with any evolutionary search method.

\subsection{Multi-Seed Analysis}
\label{sec:disc-future-seed}
% ------------------------------------------------------------

A multi-seed study is a natural extension of the present work. Running
each task with five to ten seeds and reporting mean and standard
deviation would strengthen the empirical claims. It would also allow
us to test whether the observed patterns (recurrence scaling, sparsity,
growth regimes) hold across seeds.

\subsection{Extending to Structured Inputs}
\label{sec:disc-future-structured}
% ------------------------------------------------------------

Our current encoding assumes that inputs are unstructured vectors of
bits. Many real-world tasks involve structured inputs: images have
spatial structure, sequences have temporal structure, and graphs have
relational structure. Extending \modnet{} to handle these inputs would
require incorporating structural priors into the architecture.

One approach is to use a convolutional or graph-based preprocessing
layer to extract local features, and then feed these features to a
\modnet{} network. This is analogous to the way that modern vision
models use a convolutional backbone followed by a fully connected head.
The difference is that the head in our case is a layer-free recurrent
network rather than a stack of dense layers.

Another approach is to make \modnet{} itself structure-aware, by
allowing nodes to have ``receptive fields'' that are defined by the
input structure. This is closer to the approach taken by HyperNEAT
\citep{stanley2009hypercube}, where the geometry of the substrate
encodes the structure of the input. Extending this idea to the
layer-free setting is an interesting direction.

\subsection{Combining Evolution with Gradient Descent}
\label{sec:disc-future-hybrid}
% ------------------------------------------------------------

A natural extension of \modnet{} is to combine evolutionary topology
search with gradient-based weight optimization. In this hybrid
approach, the evolutionary algorithm would discover the topology, and
gradient descent would fine-tune the weights. This is analogous to the
approach taken by \citet{real2019regularized} for image classification,
but applied to the layer-free setting.

The challenge is that the Heaviside activation is not differentiable,
so standard backpropagation cannot be applied directly. One solution is
to use a differentiable surrogate for the activation during training,
and then use the original Heaviside activation during evaluation. This
is the approach taken by spiking neural networks
\citep{neftci2019surrogate}. Another approach is to replace the
Heaviside with a sigmoid during a ``polishing'' phase, and then revert
to the Heaviside once the weights have converged.

We have explored the second approach in preliminary experiments, but
the results were not encouraging: the gradient information was not
informative enough to improve the weights significantly. A more
thorough investigation is needed to determine whether this approach can
work in practice.

\subsection{Theoretical Analysis}
\label{sec:disc-future-theory}
% ------------------------------------------------------------

Finally, a more thorough theoretical analysis of \modnet{} would be
valuable. There are several questions we would like to answer:

\begin{itemize}
  \item What is the computational power of \modnet{}? Is it Turing
    complete? Under what conditions on the number of nodes and time
    steps?
  \item What is the relationship between the evolved structure and the
    task? Can we predict the number of nodes required for a given task?
    Can we characterize the class of tasks that \modnet{} can solve
    efficiently?
  \item What is the convergence rate of the evolutionary algorithm?
    How many generations are required to reach a given level of
    accuracy? How does this depend on the task?
  \item Is the recurrence count a fundamental invariant? Are there
    tasks for which the minimum recurrence is provably $\Omega(f(n))$
    for some function $f$?
\end{itemize}

Answering these questions would require a combination of empirical
studies and theoretical analysis, and would likely require tools from
several fields (evolutionary computation, complexity theory, and
statistical learning theory).

% ------------------------------------------------------------
\section{Broader Implications}
\label{sec:disc-broader}
% ------------------------------------------------------------

Beyond the technical results, this work has broader implications for
the way we think about neural architecture design.

The dominant paradigm in modern machine learning is to design
architectures by hand, guided by intuition and empirical tuning. This
has been extremely successful, but it has also led to a proliferation
of architectures that are highly specialized for particular tasks. The
transformer, for example, is optimized for sequences; the convolutional
network is optimized for images; the recurrent network is optimized
for time series. Each architecture comes with its own set of
hyperparameters, its own training procedure, and its own failure modes.

\modnet{} suggests an alternative: a single, general-purpose
architecture that can be adapted to any task through evolution. The
same node model, the same fitness function, and the same evolutionary
algorithm solved Boolean logic, digital arithmetic, sorting networks,
and continuous regression with no task-specific modifications.

This is not to say that \modnet{} is better than specialized
architectures on their home turf---it is not, at least not at the scale
we have tested. But it does suggest that architecture design can be
automated, and that the resulting architectures need not be
specialized.

A second implication concerns the relationship between architectural
priors and task difficulty. Our comparison of layer-free and modular
variants suggests that priors that help on some tasks hurt on others.
This is a specific instance of the ``no free lunch'' theorem
\citep{wolpert1997no}: no single architecture is best for all problems,
but a single architecture combined with a powerful search algorithm
might be able to solve a wide range of problems by discovering
task-specific structures on the fly. This is the essence of
meta-learning, and \modnet{} is a step in that direction.

% ------------------------------------------------------------
\section{Summary}
\label{sec:disc-summary}
% ------------------------------------------------------------

This chapter has interpreted the results of Chapter~\ref{ch:results}
and discussed their limitations. The main conclusions are:

\begin{enumerate}
  \item \textbf{Layer-free computation is feasible for digital logic.}
    Eleven of eleven Boolean and digital-circuit tasks were solved to
    exact $100\%$ accuracy, including full adder, 2-bit adder,
    comparator, MUX, and 4-bit sorting network.

  \item \textbf{Recurrence is a structural necessity.} The number of
    self-loops scales with the temporal complexity of the task,
    ranging from one (XOR-2) to fourteen (Gaussian).

  \item \textbf{Layer-free is equal or superior to modular.} The
    layer-free variant is never worse than the modular variant on any
    task, and is decisively better on the multiplication task.

  \item \textbf{The architecture has limits.} Parity-6 was not solved
    exactly within our computational budget, and there is no reason to
    expect \modnet{} to scale indefinitely.

  \item \textbf{There is much room for improvement.} Multi-seed
    analysis, structured inputs, hybrid evolution-gradient methods,
    and theoretical analysis are all promising directions for future
    work.
\end{enumerate}

These observations motivate the concluding remarks in
Chapter~\ref{ch:conclusion}.

% ============================================================
%  Bibliography for this chapter
% ============================================================

% ============================================================
%  Chapter 7: Conclusion
%  File: chapters/07_conclusion.tex
%  Bibliography: chapters/refs_07_conclusion.bib
% ============================================================

\chapter{Conclusion}
\label{ch:conclusion}

\section{Summary of the Work}

This thesis began with a simple question: \emph{can a neural network
discover its own structure from first principles, without a designer
specifying layers, connectivity, or even the number of nodes?}

We answered this question in the affirmative, through an architecture
we call \modnet{}. The architecture has four defining properties:

\begin{enumerate}
  \item \textbf{Layer-free.} There are no predefined layers, and no
    modular grouping. Every node is a general-purpose threshold unit
    that may connect to any other node, including itself.

  \item \textbf{Outputs as ordinary nodes.} The first $n_{\text{out}}$
    nodes of the network are designated as outputs; there is no separate
    output layer.

  \item \textbf{Grown and pruned.} The network starts from a small random
    topology and grows through addition of nodes, and shrinks through
    removal of inactive or unreachable ones. Growth is bidirectional.

  \item \textbf{Gradient-free.} No gradient information is used at any
    stage. All learning is performed by a compact evolutionary algorithm
    operating on the topology and weights simultaneously.
\end{enumerate}

We evaluated \modnet{} on a suite of \textbf{sixteen} benchmark problems
spanning four broad classes: Boolean logic, digital arithmetic, sorting
networks, and continuous regression. Table~\ref{tab:conclusion-headline}
summarizes the headline results.

\begin{table}[h]
  \centering
  \small
  \caption{Headline results of \modnet{} on the sixteen benchmark
    problems. Boolean tasks are measured in exact accuracy; continuous
    tasks in root-mean-square error (RMSE).}
  \label{tab:conclusion-headline}
  \begin{tabular}{llrrr}
    \toprule
    \textbf{Task} & \textbf{Type} & \textbf{Result} & \textbf{Nodes} & \textbf{Time (s)} \\
    \midrule
    XOR-2               & Boolean    & \textbf{100.0\%} &  8 &  0.8 \\
    XOR-3               & Boolean    & \textbf{100.0\%} &  9 &  2.0 \\
    Parity-4            & Boolean    & \textbf{100.0\%} & 10 &  1.8 \\
    Parity-6            & Boolean    & 95.3\%           &  8 & 40.6 \\
    Majority-3          & Boolean    & \textbf{100.0\%} & 26 &  0.9 \\
    Majority-5          & Boolean    & \textbf{100.0\%} & 26 & 10.2 \\
    Full Adder          & Boolean    & \textbf{100.0\%} & 26 &  3.0 \\
    2-bit Adder         & Boolean    & \textbf{100.0\%} & 26 &  7.3 \\
    Comparator          & Boolean    & \textbf{100.0\%} & 27 &  3.8 \\
    MUX 4-to-1          & Boolean    & \textbf{100.0\%} & 26 &  4.8 \\
    Sort-4              & Boolean    & \textbf{100.0\%} & 26 &  8.7 \\
    \midrule
    $\sin(2\pi x)$      & Regression & RMSE $= 0.0465$  & 18 & 30.0 \\
    $\sin(4\pi x)$      & Regression & RMSE $= 0.0789$  & 13 & 25.1 \\
    $x^2 + y^2$         & Regression & RMSE $= 0.0143$  & 26 & 21.1 \\
    Gaussian            & Regression & RMSE $= 0.0275$  & 27 & 41.6 \\
    Multiplication      & Regression & RMSE $= 0.0254$  & 13 & 28.5 \\
    \bottomrule
  \end{tabular}
\end{table}

\section{What We Learned}

Beyond the raw numbers, this work produced five insights that we
believe are worth carrying forward.

\subsection{Structure Is Discoverable}

The most important lesson is that \emph{structure is discoverable}.
Given only a minimal node model and a fitness function, evolution
found:

- Minimal networks for Boolean tasks (eight nodes for XOR-2).
- Recurrent networks for temporal tasks (up to fourteen self-loops).
- Sparse networks for all tasks ($27\%$--$67\%$ density).
- Networks that solve nontrivial digital circuits exactly.

No human designer specified any of these properties. They emerged from
the interaction between growth, pruning, mutation, and selection.

\subsection{Recurrence Is a Structural Necessity, Not a Design Choice}

We observed that the number of recurrent self-loops scales with the
temporal complexity of the task:

\begin{itemize}
  \item XOR-2 (combinational): one self-loop.
  \item Parity-4 (combinational, deeper): five self-loops.
  \item Gaussian (continuous): fourteen self-loops.
  \item Multiplication (continuous, 2D): ten self-loops.
\end{itemize}

This suggests that recurrence is not a special architectural feature to
be added when needed, but a natural consequence of tasks that require
temporal integration. \modnet{} discovers this necessity without being
told.

\subsection{Sparsity Is Preferred}

Evolved networks use only a fraction of the available connections. The
pruning mechanism removes nodes that are either inactive or saturated,
and the resulting networks have densities of $27\%$ to $67\%$
depending on the task. This is consistent with the sparsity of
biological neural networks \citep{chklovskii2004wiring, cherniak1994component}.

\subsection{Modularity Is Not Necessary}

In earlier work, we studied a modular variant of \modnet{}, where
nodes were grouped into a fixed number of modules. The present work is
fully layer-free: there are no modules at all. A direct comparison
revealed that the layer-free variant is \emph{equal or superior} on
every benchmark, including a decisive win on the multiplication task
($13$ nodes versus $18$). This suggests that modular scaffolds, far
from helping, can restrict the search when the task does not require
modular structure.

\subsection{Digital Circuits Are Within Reach}

To the best of our knowledge, this is the first demonstration of a
fully layer-free evolved network that solves full adder, 2-bit adder,
comparator, 4-to-1 multiplexer, and 4-bit sorting network, all with
exact $100\%$ accuracy. This extends the scope of neuroevolution
beyond classification and control into the domain of digital logic
synthesis.

\section{Limitations}

We do not claim that \modnet{} is a universal solver. Several
limitations must be acknowledged.

\subsection{Scale}

The largest network in our experiments has $27$ nodes. This is small
by modern standards: even a modest feedforward network for a real-world
task typically has thousands of parameters. The primary reason for
this limitation is computational: the fitness evaluation is
$O(E \cdot T)$ per network, and the evolutionary algorithm performs
$O(P \cdot G)$ evaluations, so the total cost grows with the size of
the network.

\subsection{Two Failure Cases}

Two benchmark problems did not reach their success threshold:

\begin{itemize}
  \item \textbf{Parity-6}: stuck at $95.3\%$ accuracy with a network of
    eight nodes. The learning curve shows no sign of approaching exact
    accuracy, suggesting a genuine local optimum.
  \item \textbf{$\sin(4\pi x)$}: reached RMSE $= 0.0789$, close to but
    above the $0.05$ threshold. The network was still improving at the
    end of the run, suggesting that additional generations might
    suffice.
\end{itemize}

These failures are informative: they show that \modnet{} is not a
universal solver, and that certain tasks require either more resources
or a different search strategy.

\subsection{No Theoretical Guarantees}

We have not proven any theoretical properties of \modnet{}. In
particular, we have not shown that the evolutionary algorithm converges
to a global optimum, nor have we formally established Turing
completeness. The empirical results are encouraging, but they do not
constitute a proof.

\subsection{Single Seed}

Our experiments use a fixed random seed. With a different seed, the
algorithm might find different (and possibly better or worse) networks.
A multi-seed analysis would allow us to report means and standard
deviations, and to distinguish robust patterns from coincidences.

\subsection{Limited Task Diversity}

The sixteen tasks we considered all involve low-dimensional inputs
($2$--$8$ bits) and outputs ($1$--$8$ bits). This is far from the
dimensionality of real-world problems. We have not demonstrated that
\modnet{} scales to high-dimensional inputs.

\section{Future Work}

Several directions for future research emerge from this work.

\subsection{Scaling to Larger Networks}

The most obvious direction is to scale \modnet{} to larger networks.
There are several ways to approach this:

\begin{enumerate}
  \item \textbf{Parallelize the fitness evaluation.} Each network in
    the population is independent, so evaluations can be distributed
    across multiple cores or machines. With a cluster of GPUs, a
    population of $10^5$ networks could be evaluated in the time it
    takes to evaluate $10^3$ on a single machine.

  \item \textbf{Sparse encoding.} The current encoding represents each
    weight as a 32-bit float. A sparse encoding that stores only
    active connections could reduce memory and computation by an
    order of magnitude.

  \item \textbf{Advanced search algorithms.} The evolutionary algorithm
    we use is simple. More advanced methods such as CMA-ES,
    differential evolution, or Bayesian optimization might converge
    faster.
\end{enumerate}

\subsection{Hybrid Learning}

A natural extension is to combine evolutionary topology search with
gradient-based weight optimization. In this hybrid approach, the
evolutionary algorithm would discover the topology, and gradient
descent would fine-tune the weights. The challenge is that the
Heaviside activation is not differentiable; one solution is to use a
differentiable surrogate during training, as done in spiking neural
networks \citep{neftci2019surrogate}.

\subsection{Extending to Structured Inputs}

Our current encoding assumes unstructured bit vectors. Extending
\modnet{} to images, sequences, or graphs would require incorporating
structural priors into the architecture, for example through
convolutional or graph-based preprocessing.

\subsection{Theoretical Analysis}

Finally, a more thorough theoretical analysis would be valuable. Key
questions include:

\begin{itemize}
  \item Under what conditions is \modnet{} Turing complete?
  \item What is the convergence rate of the evolutionary algorithm?
  \item Can we characterize the class of tasks that \modnet{} can
    solve efficiently?
\end{itemize}

\section{Broader Implications}

Beyond the technical results, this work has broader implications for
the way we think about neural architecture design.

The dominant paradigm in modern machine learning is to design
architectures by hand, guided by intuition and empirical tuning. This
has been remarkably successful, but it has also led to a proliferation
of architectures that are highly specialized for particular tasks. The
transformer, for example, is optimized for sequences; the
convolutional network is optimized for images; the recurrent network
is optimized for time series. Each architecture comes with its own set
of hyperparameters, its own training procedure, and its own failure
modes.

\modnet{} suggests an alternative: a single, general-purpose
architecture that can be adapted to any task through evolution. The
same node model, the same fitness function, and the same evolutionary
algorithm solved Boolean logic, digital arithmetic, sorting, and
continuous regression with no task-specific modifications.

This is not to say that \modnet{} is better than specialized
architectures on their home turf---it is not, at least not at the
scale we have tested. But it does suggest that architecture design
can be automated, and that the resulting architectures need not be
specialized. This is the essence of meta-learning, and \modnet{} is a
step in that direction.

\section{Final Remarks}

We began this thesis with the observation that biological brains have
no layers, no fixed topology, and no global error signal. They
discover their structure through developmental and evolutionary
processes. \modnet{} is a modest attempt to capture these properties
in a minimal computational model.

Our results suggest that this approach is viable. A single
layer-free architecture, driven by a compact evolutionary process,
solved fourteen of sixteen tasks spanning Boolean logic, digital
arithmetic, sorting networks, and continuous regression, using no more
than $27$ nodes per network. Recurrence emerged spontaneously and
scaled with task complexity. Sparsity emerged without any regularizer.
And a direct comparison with a modular variant showed that removing
structure entirely can \emph{improve} performance.

There is much work still to be done. The architecture has limits---two
of our sixteen tasks were not fully solved---and scaling remains an
open challenge. But the direction is promising. If structure can be
discovered rather than designed, then the field of neural architecture
may be entering a new phase, one in which the designer's role shifts
from \emph{building} networks to \emph{specifying objectives}.

We close with a reflection. The history of machine learning is a
history of \emph{removing} structure from our models. Early models were
heavily hand-crafted. Deep learning removed the feature engineering.
Neural architecture search removed the architecture design. \modnet{}
takes this trend one step further: it removes even the layered
structure itself, and asks the network to discover it.

Perhaps the next step is to remove ourselves entirely from the design
loop. If a network can discover its own structure, perhaps it can
also discover its own objective. That question is beyond the scope of
this thesis, but the tools developed here---self-growing layer-free
networks with complete event logging and reproducible benchmarks---may
serve as a foundation for exploring it.

% ============================================================
%  Bibliography for this chapter
% ============================================================


% ============================================================
%                       APPENDICES
% ============================================================
\appendix
\chapter{Hyperparameter Summary}
\label{app:hyperparameters}

Table~\ref{tab:app-all-hyperparameters} lists all hyperparameters
used in the experiments reported in this thesis.

\begin{table}[h]
  \centering
  \small
  \caption{Complete list of hyperparameters for \modnet{}.}
  \label{tab:app-all-hyperparameters}
  \begin{tabular}{lll}
    \toprule
    \textbf{Parameter} & \textbf{Description} & \textbf{Value} \\
    \midrule
    $P$                & Population size              & 300 \\
    $G$                & Maximum generations          & 400--1500 \\
    $p_{\text{cross}}$ & Crossover rate               & 0.30 \\
    $p_{\text{fresh}}$ & Fresh random rate            & 0.05 \\
    $\sigma_{\text{base}}$ & Base mutation scale      & 0.30 \\
    $\sigma_{\max}$    & Maximum mutation scale       & 1.50 \\
    $\gamma$           & Mutation growth factor       & 1.25 \\
    $g_{\text{thresh}}$& Plateau threshold            & 25 \\
    $g_{\text{birth}}$ & Birth interval               & 100 \\
    $g_{\text{prune}}$ & Prune interval               & 20 \\
    $\epsilon_{\text{prune}}$ & Prune tolerance       & $-0.001$ \\
    $w_{\text{thresh}}$& Weight pruning threshold     & 0.04 \\
    $\alpha_{\text{lo}}$& Lower activity threshold    & 0.02 \\
    $\alpha_{\text{hi}}$& Upper activity threshold    & 0.98 \\
    $n_{\text{hidden,min}}$& Minimum hidden nodes    & 2 \\
    $n_{\max}$         & Maximum nodes                & 60 \\
    \bottomrule
  \end{tabular}
\end{table}


\chapter{Network Snapshots}
\label{app:snapshots}

This appendix presents network snapshots at selected generations for
each benchmark problem, using the spring layout of
Section~\ref{sec:exp-viz}. Full-resolution PNGs are in the
open-source repository at \url{https://github.com/mablue92/modnet}.


\chapter{Algorithm Details}
\label{app:algorithm}

\begin{algorithm}[h]
  \caption{ModNet Evolution (full detail)}
  \label{alg:app-evolution}
  \begin{algorithmic}[1]
    \State Initialize population $\mathcal{P}$ with $P$ random networks
    \State Evaluate fitness for all networks
    \For{$g = 1, 2, \ldots, G$}
      \State Sort $\mathcal{P}$ by fitness
      \State $\mathcal{E} \gets$ top $\lfloor P/5 \rfloor$ networks
      \State Apply crossover, mutation, and fresh injection
      \If{$g \bmod g_{\text{birth}} = 0$}
        \State Apply \textsc{Birth} to top 3 networks
      \EndIf
      \If{$g \bmod g_{\text{prune}} = 0$}
        \State Apply \textsc{Prune} to elite networks
      \EndIf
    \EndFor
    \State \Return best network
  \end{algorithmic}
\end{algorithm}




% -------------------------------------------------------------
%  Consolidated bibliography (replaces per-chapter bibs)
% -------------------------------------------------------------
\bibliographystyle{plainnat}
\bibliography{refs_all}

\end{document}